\documentclass[11pt,a4paper]{article}
\usepackage[utf8]{inputenc}
\usepackage[T1]{fontenc}
\usepackage{geometry}
\geometry{margin=2.5cm}
\usepackage{longtable}
\usepackage{booktabs}
\usepackage{array}
\usepackage{colortbl}
\usepackage{xcolor}
\usepackage{graphicx}
\usepackage{hyperref}
\usepackage{enumitem}
\usepackage{titlesec}
\usepackage{fancyhdr}
\usepackage{tabularx}
\usepackage{multirow}
\usepackage{amsmath}
\usepackage{amssymb}
\usepackage{float}
\usepackage{pdfpages}

\definecolor{uteqblue}{RGB}{0,51,102}

\hypersetup{
  colorlinks=true,
  linkcolor=uteqblue,
  citecolor=uteqblue,
  urlcolor=blue,
  pdftitle={Especificación de Requisitos de Software (ERS/SRS) v2.0 - CarniCore},
  pdfauthor={Equipo CarniCore - UTEQ ISR-401}
}

\pagestyle{fancy}
\fancyhf{}
\fancyhead[L]{UTEQ -- ISR-401 / 20303}
\fancyhead[R]{ERS/SRS -- CarniCore -- Entrega 4 (2B/Defensa Final)}
\fancyfoot[C]{\thepage}

\begin{document}

%==========================================================
% PORTADA
%==========================================================
\begin{titlepage}
    \centering

    \textbf{\large UNIVERSIDAD TÉCNICA ESTATAL DE QUEVEDO}\\[0.3cm]
    \textbf{FACULTAD DE CIENCIAS DE LA COMPUTACIÓN Y DISEÑO DIGITAL}\\[0.3cm]
    \textbf{CARRERA DE INGENIERÍA DE SOFTWARE (REDISEÑO)}\\[0.1cm]

    \includegraphics[width=4cm]{figures/logo.png}\\[0.1cm]


    \rule{\textwidth}{1pt}\\[0.4cm]
    {\Large \textbf{ESPECIFICACIÓN DE REQUISITOS DE SOFTWARE (ERS/SRS)}}\\[0.2cm]
    {\large \textbf{Sistema CarniCore}}\\[0.2cm]
    {\normalsize Sistema de Trazabilidad, Pesaje Inteligente e Inventario\\
    Distribuidora de Cárnicos Pucayacu -- La Maná, Cotopaxi}\\[0.2cm]
    {\normalsize \textbf{Entrega 4 (2B / Defensa Final): Versión Definitiva v2.4.0}}\\[0.2cm]

    \rule{\textwidth}{1pt}\\[0.3cm]

    \textbf{INTEGRANTES:}\\[0.3cm]
    \begin{tabular}{ll}
    CASTRO BAJA\~NA ARIEL OMAR     & Cédula: 2350262305 \\[0.2cm]
    CRESPO ESPINOZA KLEBER OBED    & Cédula: 1315380244 \\[0.2cm]
    GAMARRA ARAUJO EDHU XAVIER     & Cédula: 1208370633 \\[0.2cm]
    PÉREZ RUIZ CARLOS ANDRÉS       & Cédula: 0955713136 \\[0.2cm]
    QUINTERO GENDE ERICK JAHIR     & Cédula: 1250575527 \\
    \end{tabular}\\[0.5cm]

    \textbf{ASIGNATURA:} Ingeniería de Requerimientos [20303] -- ISR-401\\[0.3cm]
    \textbf{CURSO:} 4to ``A'' Software\\[0.3cm]
    \textbf{DOCENTE:} PhD. Guerrero Ulloa Gleiston Cicerón\\[0.3cm]
    \textbf{PERÍODO:} 2026--2027 PPA\\[0.3cm]
    \textbf{FECHA:} Semana 20 -- septiembre 14 del 2026\\[0.3cm]
    \textbf{REPOSITORIO:}\\[0.2cm]
    \url{https://github.com/gleiston-guerrero/CarniCore_Proyecto_ISR401.git}\\[0.4cm]
   
    \textbf{DEPÓSITO FAIR (ZENODO):}\\[0.2cm]
    DOI: \url{https://doi.org/10.5281/zenodo.22225854}\\[0.2cm]
    \textbf{REGISTRO OSF:}\\[0.2cm]
    \url{https://osf.io/yp7t3}\\[0.4cm]
    \vfill
\end{titlepage}

%==========================================================
% HISTORIAL DE VERSIONES
%==========================================================
\newpage
\section*{Historial de Versiones}
\begin{longtable}{|p{1.5cm}|p{2.2cm}|p{3cm}|p{8cm}|}
\caption{Historial de versiones del documento ERS/SRS -CarniCore}\\
\hline
\textbf{Versión} & \textbf{Fecha} & \textbf{Autor(es)} & \textbf{Descripción del cambio} \\
\hline
1.0 & 31/05/2026 & Equipo CarniCore & Entrega 1A: planificación del proyecto de IR, introducción y descripción general preliminar. \\
\hline
1.5 & 15/06/2026 & Equipo CarniCore & Correcciones sobre la 1A: alcance delimitado, exclusión explícita del módulo de facturación electrónica del SRI. \\
\hline
2.0 & 12/07/2026 & Equipo CarniCore & Entrega 1B: interfaces externas con mockups, 15 RF, 6 RNF, restricciones de diseño, modelado UML parcial, MoSCoW y matriz de trazabilidad parcial. \\
\hline
3.0 & 02/08/2026 & Equipo CarniCore & Entrega 3 (2A): ERS completo (25 RF, 15 RNF) con 8 atributos y EV-XX; modelado UML completo; priorización MoSCoW+Kano+WSJF; matriz de trazabilidad extendida (40 filas); componente empírico y puente hacia publicación; registro en OSF. \\
\hline
\textbf{4.0 / v2.0} & \textbf{21/08/2026} & \textbf{Equipo CarniCore} & \textbf{Entrega 4 (2B) -- Versión definitiva:} incorporación de RF-26 y RF-27 (CCB aprobados); diagrama i* SR completado; requisitos de IA (IA-01 e IA-02); auditoría de calidad con seis métricas ISO/IEC 25010:2023; trazabilidad end-to-end (60 filas: 42 base + 18 IA; ampliada a 84 en el examen suspenso del 16/09/2026 con 24 filas de flujos alternativos y excepciones); re-inspección Fagan PE5 con 4 defectos adicionales corregidos; retrospectiva del equipo; conclusiones integradoras; enlace al DOI de Zenodo en portada. \\
\hline
\textbf{4.1 / v2.1} & \textbf{16/09/2026} & \textbf{Equipo CarniCore} & \textbf{Corrección de cierre (examen suspenso):} se especificó al menos un flujo alternativo y una excepción, con condición de disparo y poscondición propias, en los 12 casos de uso (antes solo dos tenían flujo alternativo explícito y cuatro excepción); se resolvieron cinco contradicciones entre precondición y flujo detectadas en revisión (CU-01, CU-02, CU-08, CU-11, CU-12); se incorporó la georreferencia de RF-17 como flujo alternativo propio de CU-01; se crearon HU-05 y HU-12 para cubrir los flujos de CU-05 y CU-11, que no tenían historia de usuario asociada; matriz de trazabilidad ampliada de 60 a 84 filas (24 filas nuevas de flujos alternativos y excepciones, columna Flujo añadida) y sincronizada con \texttt{04\_Trazabilidad/Matriz\_Trazabilidad.csv} (90 filas totales, incluidas 6 filas de RNF/CU sin correspondencia directa en esta tabla). \\
\hline
\end{longtable}

\renewcommand{\contentsname}{\'{I}ndice}
\renewcommand{\listtablename}{\'{I}ndice de Tablas}
\renewcommand{\listfigurename}{\'{I}ndice de Figuras}

\newpage
\tableofcontents
\newpage
\listoftables
\newpage
\listoffigures
\newpage

%==========================================================
\section{Introducción}
%==========================================================
\subsection{Propósito}

El presente documento constituye la Especificación de Requisitos de Software (ERS/SRS) definitiva del sistema CarniCore, elaborada siguiendo la estructura y terminología del estándar ISO/IEC/ IEEE~29148:2018 \cite{iso29148}. Esta versión final (Entrega 4 -- 2B) cierra el ciclo abierto en la Entrega 1A: incorpora la totalidad de las observaciones docentes de la Entrega 3 (2A), los requisitos adicionales aprobados por el CCB (RF-26 y RF-27), el diagrama i*~SR completo, la especificación de dos componentes de Inteligencia Artificial (IA-01 e IA-02), la auditoría de calidad con seis métricas y la trazabilidad consolidada de 84 filas (60 base/IA + 24 de flujos alternativos y excepciones, añadidas en el cierre del examen suspenso del 16/09/2026).

Este documento está dirigido a tres audiencias: (1) el equipo de desarrollo, que lo usará como insumo definitivo para el diseño arquitectónico y la construcción del MVP; (2) el personal docente y evaluador de la asignatura Ingeniería de Requerimientos [20303], que verificará el cumplimiento de los criterios C1--C13 de la rúbrica de la Entrega 4 (2B); y (3) la propietaria de la distribuidora y su equipo administrativo, en calidad de cliente y validador funcional final.

\subsection{Alcance}

El software se denomina \textbf{CarniCore} y digitaliza y da trazabilidad a los procesos operativos críticos de una distribuidora de productos cárnicos (res, cerdo, pollo y embutidos) de Pucayacu, La Maná, Cotopaxi, actualmente gestionados de forma manual.

\paragraph{El sistema SÍ hará:}
\begin{itemize}[nosep]
\item Registrar el ingreso de lotes y animales (cerdo y res con número de arete) provenientes de proveedores certificados, incluida la guía de origen y movilización.
\item Registrar el pesaje de productos y calcular automáticamente el costo por libra, emitiendo un comprobante de pesaje.
\item Registrar el proceso de despiece y clasificación de cortes.
\item Gestionar el inventario en cámaras frigoríficas y controlar los días de vida útil según el tipo de producto.
\item Registrar bajas de producto con evidencia fotográfica y aprobación del administrador.
\item Permitir la consulta de trazabilidad completa desde el proveedor hasta la venta final, incluyendo vista pública restringida mediante QR con campos delimitados (RF-26).
\item Generar reportes gerenciales filtrables por proveedor (RF-27) y por período.
\item Gestionar usuarios con distintos roles y operar en modo offline, sincronizando posteriormente.
\item Proporcionar predicción de demanda y detección de anomalías de cadena de frío (componentes IA-01 e IA-02).
\end{itemize}

\paragraph{El sistema NO hará (exclusiones explícitas):}
\begin{itemize}[nosep]
\item No sustituirá al sistema de facturación electrónica del SRI.
\item No emitirá guías de movilización ante Agrocalidad.
\item No incluirá comercio electrónico ni venta en línea al consumidor final.
\item No controlará de forma autónoma sensores IoT de temperatura en esta fase.
\end{itemize}

\subsection{Glosario de términos, acrónimos y abreviaturas}

\begin{longtable}{|p{3.5cm}|p{11cm}|}
\caption{Glosario de términos, acrónimos y abreviaturas}\\
\hline
\textbf{Término} & \textbf{Definición} \\
\hline
ERS/SRS & Especificación de Requisitos de Software / \textit{Software Requirements Specification}. \\
RF / RNF & Requisito Funcional / Requisito No Funcional. \\
CU & Caso de Uso. \\
HU & Historia de Usuario (formato Connextra). \\
CA & Criterio de Aceptación (formato Gherkin: Dado--Cuando--Entonces). \\
EV-\textit{XX} & Identificador de Evidencia que respalda un requisito. \\
RD & Restricción de Diseño. \\
MU-\textit{XX} & Identificador de mockup. \\
Canal & Cuerpo del animal faenado, sin cabeza ni vísceras, listo para despiece. \\
Faenamiento & Proceso de sacrificio y procesamiento del animal en la planta. \\
Merma de faenamiento & Diferencia entre el peso vivo del animal y el peso de la canal obtenida tras el faenamiento. \\
Merma de despiece & Diferencia entre el peso de la canal y la suma de los cortes registrados; margen configurable (defecto 3\%, rango 1--8\%). \\
Guía de origen & Documento oficial de Agrocalidad que certifica la procedencia y el número de arete de los animales movilizados. \\
Arete & Identificador físico único colocado en el animal para su trazabilidad ante Agrocalidad. \\
Agrocalidad & Agencia de Regulación y Control Fito y Zoosanitario del Ecuador. \\
MoSCoW & Técnica de priorización: \textit{Must have, Should have, Could have, Won't have}. \\
Kano & Modelo de clasificación de características según satisfacción del cliente. \\
WSJF & \textit{Weighted Shortest Job First}, técnica SAFe para priorizar por costo de retraso. \\
CCB & Comité de Control de Cambios (\textit{Change Control Board}). \\
RFC & Solicitud de Cambio (\textit{Request for Change}). \\
DOI & Identificador de Objeto Digital; identifica el paquete de datos depositado en Zenodo. \\
LOPDP & Ley Orgánica de Protección de Datos Personales del Ecuador \cite{lopd2021}. \\
IA & Inteligencia Artificial. \\
MVP & Producto Mínimo Viable. \\
ISO/IEC 25010 & Norma de calidad de producto de software (SQuaRE) \cite{iso25010}. \\
i* SD & Diagrama de Dependencia Estratégica del marco \textit{i-star}. \\
i* SR & Diagrama de Razonamiento Estratégico del marco \textit{i-star}. \\
\hline
\end{longtable}

\subsection{Referencias}

La bibliografía completa se gestiona en \texttt{referencias.bib} con al menos 40 entradas verificadas, incluyendo las normas base \cite{iso29148,iso25010}, fundamentos de IR \cite{pohl2010,swebok2024}, principios FAIR \cite{wilkinson2016fair}, literatura sobre IA generativa en IR \cite{cheng2026genai,arora2024genai,vogelsang2024llm}, explicabilidad como RNF \cite{chazette2020explainability,chazette2021exploring,chazette2022explainable}, la LOPDP \cite{lopd2021}, guías de estudios empíricos \cite{molleri2020checklist,hennink2017saturation,runeson2009casestudy,wohlin2012experimentation,nosek2018preregistration,cohen1960kappa,fleiss1971kappa,ferrari2016ambiguity,fischbach2023acceptance,kitchenham2007slr} y trazabilidad cárnica \cite{chen2021blockchain,beck2021blockchainmeat}.

\subsection{Visión general del documento}

El documento se organiza en: Sección~2 (descripción general, incluyendo modelos i*~SD y SR); Sección~3 (requisitos específicos completos: RF-01 a RF-27, RNF-01 a RNF-15, requisitos legales, HU y restricciones de diseño); Sección~4 (modelado UML completo); Sección~5 (priorización y matriz de trazabilidad de 84 filas); Sección~6 (MVP); Sección~7 (componente empírico y protocolo registrado en OSF); Sección~8 (requisitos de IA: IA-01 e IA-02); Sección~9 (auditoría de calidad con seis métricas); y apéndices con plan del proyecto, evidencias, retrospectiva y declaración de uso de IA.

%==========================================================
\section{Descripción General}
%==========================================================
\subsection{Perspectiva del producto}

CarniCore es un sistema nuevo; la distribuidora no cuenta actualmente con ningún sistema informático de gestión interna. El sistema es independiente y no reemplaza a ningún software existente, aunque referenciará el número de factura generado por el sistema contable externo (SRI).

\begin{figure}[H]
    \centering
    \includegraphics[width=\textwidth,keepaspectratio]{figures/Context_Diagram.png}
    \caption{Diagrama de Contexto del Sistema CarniCore -- relaciones externas e interfaces}
    \label{fig:context}
\end{figure}

\subsection{Funciones del producto}
\begin{enumerate}[nosep]
\item Gestión de proveedores y guías de origen/movilización.
\item Registro de ingreso de lotes y animales.
\item Pesaje inteligente con cálculo automático de costos.
\item Registro de despiece y clasificación de cortes.
\item Gestión de cámaras frigoríficas e inventario por sección.
\item Control de vida útil y alertas de vencimiento.
\item Registro de bajas de producto con evidencia y supervisión.
\item Consulta de trazabilidad extremo a extremo, incluida vista pública vía QR con campos delimitados (RF-26).
\item Registro de ventas, salidas y devoluciones de producto.
\item Generación de reportes gerenciales filtrables por proveedor (RF-27) y por período.
\item Gestión de usuarios y roles.
\item Operación en modo de contingencia offline con sincronización posterior.
\item Predicción de demanda por tipo de corte (IA-01).
\item Detección de anomalías de cadena de frío (IA-02).
\end{enumerate}

\subsection{Mapa de interesados (stakeholders) y matriz poder/interés}

\begin{longtable}{|p{3.2cm}|p{2.5cm}|p{2cm}|p{2cm}|p{5cm}|}
\caption{Mapa de interesados y matriz poder/interés}\\
\hline
\textbf{Interesado} & \textbf{Tipo} & \textbf{Poder} & \textbf{Interés} & \textbf{Estrategia de gestión} \\
\hline
Propietaria (Gisela E. Ibáñez Medina) & Interno / decisor & Alto & Alto & Gestionar de cerca; validación semanal de requisitos. \\
\hline
Administrador general & Administrativo interno & Alto & Alto & Gestionar de cerca; aprueba bajas e inventario. \\
\hline
Encargado de bodega & Operativo interno & Medio & Alto & Mantener informado; usuario intensivo del sistema. \\
\hline
Carnicero & Operativo interno & Bajo & Alto & Mantener informado; usuario del módulo de despiece. \\
\hline
Cliente / Sucursal & Externo & Bajo & Medio & Monitorear; consulta de trazabilidad pública. \\
\hline
Agrocalidad (ente de control) & Externo regulador & Alto & Bajo & Mantener satisfecho; exportación de evidencia de trazabilidad. \\
\hline
Proveedor / Granja & Externo & Bajo & Bajo & Monitorear; provee guías de origen. \\
\hline
Docente supervisor ISR-401 & Interno académico & Alto & Alto & Gestionar de cerca; gatekeeper de evaluación. \\
\hline
\end{longtable}

\subsection{Modelado organizacional i* -- Diagramas SD y SR}

\subsubsection{Diagrama de Dependencia Estratégica (i* SD)}

El diagrama i*~SD muestra las dependencias entre actores del sistema CarniCore:

\begin{itemize}[nosep]
\item \textbf{Propietaria $\leftrightarrow$ CarniCore:} dependencia de dato: reportes gerenciales confiables (meta blanda) desde CarniCore hacia la Propietaria; dependencia de recurso: datos de ventas/compras/pérdidas desde la Propietaria hacia CarniCore.
\item \textbf{Administrador General $\rightarrow$ CarniCore:} dependencia de tarea: aprobar bajas de producto.
\item \textbf{CarniCore $\rightarrow$ Administrador General:} dependencia de meta: alerta de vida útil procesada.
\item \textbf{Encargado de Bodega $\rightarrow$ CarniCore:} dependencia de tarea: registrar ingreso y pesaje.
\item \textbf{Carnicero $\rightarrow$ CarniCore:} dependencia de recurso (peso de cortes): registrar despiece.
\item \textbf{Agrocalidad $\rightarrow$ CarniCore:} dependencia de dato: exportación de evidencia de trazabilidad.
\item \textbf{Proveedor/Granja $\rightarrow$ CarniCore:} dependencia de recurso: guías de origen y movilización.
\end{itemize}

\begin{figure}[H]
    \centering
    \includegraphics[width=\textwidth,keepaspectratio]{figures/Strategic_Dependency_Diagram.png}
    \caption{Diagrama de Dependencia Estratégica (i* SD) del sistema CarniCore}
    \label{fig:istar_sd}
\end{figure}

\subsubsection{Diagrama de Razonamiento Estratégico (i* SR) -- Propietaria}

\textbf{Nota para inserción:} A continuación se debe insertar el diagrama i*~SR del actor Propietaria, que muestra el razonamiento interno para alcanzar sus metas blandas (\textit{softgoals}): \textit{Rentabilidad del negocio}, \textit{Cumplimiento regulatorio ante Agrocalidad} y \textit{Control de pérdidas por vencimiento}. El SR descompone las metas en tareas (registrar ventas, generar reportes, consultar trazabilidad) y recursos (datos de pesaje, datos de inventario), con contribuciones positivas (+) de los módulos de CarniCore a cada softgoal.

\begin{figure}[H]
    \centering
    \includegraphics[width=\textwidth,keepaspectratio]{figures/istar_SR_Propietaria_EN.png}
    \caption{Diagrama de Razonamiento Estratégico (i* SR) -- Actor: Propietaria}
    \label{fig:istar_sr_propietaria}
\end{figure}


\subsubsection{Diagrama de Razonamiento Estratégico (i* SR) -- Administrador General}

\begin{figure}[H]
    \centering
    \includegraphics[width=\textwidth,keepaspectratio]{figures/istar_SR_AdminGeneral_EN.png}
    \caption{Diagrama de Razonamiento Estratégico (i* SR) -- Actor: Administrador General}
    \label{fig:istar_sr_admingeneral}
\end{figure}


\subsection{Características de los usuarios}

\begin{longtable}{|p{2.8cm}|p{5cm}|p{2cm}|p{2.5cm}|}
\caption{Características de los usuarios del sistema}\\
\hline
\textbf{Tipo de usuario} & \textbf{Descripción} & \textbf{Nivel técnico} & \textbf{Frecuencia} \\
\hline
Propietaria & Lidera el equipo, coordina compras y ventas; requiere reportes gerenciales resumidos. & Básico & Diaria \\
\hline
Administrador general & Supervisa inventario, autoriza bajas, revisa diferencias. & Medio & Diaria \\
\hline
Encargado de bodega & Registra ingreso, pesaje y distribución de productos. & Básico & Diaria (múltiples veces) \\
\hline
Carnicero & Ejecuta el despiece; consulta pedidos asignados. & Básico & Diaria \\
\hline
Cliente / Sucursal & Realiza pedidos y consulta el historial de un producto. & Básico & Ocasional \\
\hline
\end{longtable}

\subsection{Entorno operativo}

Estaciones de trabajo (PC o tablet) en el punto de pesaje; dispositivo móvil para el encargado de bodega; navegador web moderno; conectividad intermitente en la planta de La Maná, por lo que se requiere modo de contingencia offline (RF-21, RNF-02, RNF-14).

\subsection{Restricciones}
\begin{itemize}[nosep]
\item \textbf{Regulatorias:} compatibilidad con formatos de guía de origen de Agrocalidad y con registros sanitarios firmados por el médico veterinario.
\item \textbf{Tecnológicas:} módulo de pesaje integrable a futuro con básculas electrónicas (salida serial/USB).
\item \textbf{Legales:} cumplimiento de la LOPDP \cite{lopd2021} para todo dato personal tratado por el sistema.
\end{itemize}

\subsection{Suposiciones y dependencias}
\begin{itemize}[nosep]
\item Se asume conexión a internet estable al menos en el punto central de bodega para sincronización.
\item Se asume que el personal recibirá capacitación básica de al menos 4 horas antes del despliegue (RNF-04).
\item Se asume que los formatos de guía de origen de Agrocalidad no cambiarán durante el desarrollo.
\end{itemize}

%==========================================================
\section{Requisitos Específicos Completos}
%==========================================================

\subsection{Requisitos de interfaces externas}

\subsubsection{Interfaces de usuario}

Se definen 20 mockups de alta fidelidad (MU-01 a MU-20) trazados a los RF correspondientes.

\begin{figure}[H]
    \centering
    \includegraphics[width=0.85\textwidth,keepaspectratio]{figures/MU-01_login.jpg}
    \caption{MU-01 -- Interfaz de Login del Sistema}
\end{figure}

\begin{figure}[H]
    \centering
    \includegraphics[width=0.85\textwidth,keepaspectratio]{figures/MU-02_dashboard.jpg}
    \caption{MU-02 -- Dashboard del Sistema}
\end{figure}

\begin{figure}[H]
    \centering
    \includegraphics[width=0.85\textwidth,keepaspectratio]{figures/MU-03_inicio.jpg}
    \caption{MU-03 -- Pantalla de Inicio}
\end{figure}

\begin{figure}[H]
    \centering
    \includegraphics[width=0.85\textwidth,keepaspectratio]{figures/MU-04_proveedores.jpg}
    \caption{MU-04 -- Módulo de Proveedores (RF-01, RF-17)}
\end{figure}

\begin{figure}[H]
    \centering
    \includegraphics[width=0.85\textwidth,keepaspectratio]{figures/MU-05_guias_origen.jpg}
    \caption{MU-05 -- Módulo de Guías de Origen (RF-02)}
\end{figure}

\begin{figure}[H]
    \centering
    \includegraphics[width=0.85\textwidth,keepaspectratio]{figures/MU-06_ingreso_lotes.jpg}
    \caption{MU-06 -- Módulo de Ingreso de Lotes (RF-03)}
\end{figure}

\begin{figure}[H]
    \centering
    \includegraphics[width=0.85\textwidth,keepaspectratio]{figures/MU-07_pesaje_inteligente.jpg}
    \caption{MU-07 -- Módulo de Pesaje Inteligente (RF-04, RF-05, RF-19)}
\end{figure}

\begin{figure}[H]
    \centering
    \includegraphics[width=0.85\textwidth,keepaspectratio]{figures/MU-08_despiece.jpg}
    \caption{MU-08 -- Módulo de Despiece (RF-06)}
\end{figure}

\begin{figure}[H]
    \centering
    \includegraphics[width=0.85\textwidth,keepaspectratio]{figures/MU-09_camaras_frigorificas.jpg}
    \caption{MU-09 -- Módulo de Cámaras Frigoríficas (RF-07)}
\end{figure}

\begin{figure}[H]
    \centering
    \includegraphics[width=0.85\textwidth,keepaspectratio]{figures/MU-10_inventario.jpg}
    \caption{MU-10 -- Módulo de Inventario (RF-07, RF-16, RF-24)}
\end{figure}

\begin{figure}[H]
    \centering
    \includegraphics[width=0.85\textwidth,keepaspectratio]{figures/MU-11_vida_util.jpg}
    \caption{MU-11 -- Módulo de Vida Útil (RF-08)}
\end{figure}

\begin{figure}[H]
    \centering
    \includegraphics[width=0.85\textwidth,keepaspectratio]{figures/MU-12_alertas.jpg}
    \caption{MU-12 -- Módulo de Alertas (RF-08)}
\end{figure}

\begin{figure}[H]
    \centering
    \includegraphics[width=0.85\textwidth,keepaspectratio]{figures/MU-13_baja_productos.jpg}
    \caption{MU-13 -- Módulo de Baja de Productos (RF-09)}
\end{figure}

\begin{figure}[H]
    \centering
    \includegraphics[width=0.85\textwidth,keepaspectratio]{figures/MU-14_ventas.jpg}
    \caption{MU-14 -- Módulo de Ventas (RF-10, RF-18)}
\end{figure}

\begin{figure}[H]
    \centering
    \includegraphics[width=0.85\textwidth,keepaspectratio]{figures/MU-15_movimientos_inventario.jpg}
    \caption{MU-15 -- Módulo de Movimientos de Inventario (RF-16)}
\end{figure}

\begin{figure}[H]
    \centering
    \includegraphics[width=0.85\textwidth,keepaspectratio]{figures/MU-16_trazabilidad.jpg}
    \caption{MU-16 -- Módulo de Trazabilidad (RF-11, RF-23, RF-25, RF-26)}
\end{figure}

\begin{figure}[H]
    \centering
    \includegraphics[width=0.85\textwidth,keepaspectratio]{figures/MU-17_conteo_inventario.jpg}
    \caption{MU-17 -- Módulo de Conteo de Inventario (RF-12)}
\end{figure}

\begin{figure}[H]
    \centering
    \includegraphics[width=0.85\textwidth,keepaspectratio]{figures/MU-18_reportes.jpg}
    \caption{MU-18 -- Módulo de Reportes (RF-13, RF-14, RF-27)}
\end{figure}

\begin{figure}[H]
    \centering
    \includegraphics[width=0.85\textwidth,keepaspectratio]{figures/MU-19_comprobante_pesaje.jpg}
    \caption{MU-19 -- Comprobante de Pesaje (RF-05)}
\end{figure}

\begin{figure}[H]
    \centering
    \includegraphics[width=0.85\textwidth,keepaspectratio]{figures/MU-20_roles.jpg}
    \caption{MU-20 -- Módulo de Roles y Usuarios (RF-15)}
\end{figure}

\subsubsection{Interfaces de hardware}

Estación de trabajo (PC o tablet) en el punto de pesaje con báscula Ohaus Scout STX Series (protocolo ASCII, puerto USB/serial); impresora térmica de tickets (trabajo futuro); dispositivo móvil o tablet para el encargado de bodega.

\subsubsection{Interfaces de software}

API REST interna con JSON y JWT; módulo de exportación PDF compatible con auditorías de Agrocalidad; módulo de sincronización offline (RF-21, sync $<$5 min).

\subsubsection{Interfaces de comunicación}

Comunicación cliente-servidor mediante HTTPS; notificaciones internas (alertas de vencimiento) mediante la propia interfaz web/móvil.

%----------------------------------------------------------
\subsection{Requisitos funcionales}
%----------------------------------------------------------

Se documentan \textbf{27 requisitos funcionales} con los ocho atributos exigidos (ID, Nombre, Descripción, Fuente, Prioridad, Estabilidad, Criterio de aceptación, Dependencias/EV). RF-01 a RF-25 provienen de las entregas anteriores; RF-26 y RF-27 fueron aprobados por el CCB en la Entrega 4 (2B).

\newcommand{\rfitem}[8]{%
\begin{longtable}{|p{2.4cm}|p{11.2cm}|}
\caption[RF-#1: #2]{Ficha de requisito funcional #1 -- #2}\\
\hline
\multicolumn{2}{|l|}{\textbf{#1 -- #2}} \\
\hline
ID & #1 \\
\hline
Nombre & #2 \\
\hline
Descripción & #3 \\
\hline
Fuente & #4 \\
\hline
Prioridad & #5 \\
\hline
Estabilidad & #6 \\
\hline
Criterio de aceptación & #7 \\
\hline
Dependencias / EV & #8 \\
\hline
\end{longtable}
\vspace{-0.4cm}
}

\rfitem{RF-01}{Registrar proveedor}{El sistema deberá permitir registrar un proveedor de animales o derivados cárnicos, almacenando su RUC (validado con el algoritmo de verificación del SRI, dígito verificador módulo 11), nombre y estado de regularización.}{Propietaria -- Bloque 2, Pregunta 4}{Must}{Estable}{El sistema no permite registrar un proveedor sin un número de RUC de 13 dígitos válido según el algoritmo de dígito verificador módulo 11 del SRI.}{Dependencias: RF-02. EV-01.}

\rfitem{RF-02}{Registrar guía de origen}{El sistema deberá permitir registrar la guía de origen entregada por el proveedor, incluyendo número de guía, cantidad de canales y confirmación de firma del médico veterinario de planta (RFC-03 aprobado: firma veterinaria digitalizada).}{Propietaria -- Bloque 2, Pregunta 5}{Must}{Estable}{El sistema rechaza el registro de un lote si no existe una guía de origen con firma veterinaria asociada.}{Dependencias: RF-01, RF-03. EV-02.}

\rfitem{RF-03}{Registrar ingreso de lote o animal}{El sistema deberá permitir registrar el ingreso de un lote (aves) o de animales individuales (cerdo, res) identificados por número de arete, incluyendo fecha, cantidad y tipo de producto.}{Propietaria -- Bloque 3, Pregunta 10}{Must}{Estable}{Para cerdo y res, el sistema exige un número de arete único por animal; para aves, exige únicamente la cantidad ingresada.}{Dependencias: RF-02. EV-03.}

\rfitem{RF-04}{Registrar pesaje de producto}{El sistema deberá permitir registrar el peso en libras de un producto y calcular automáticamente el costo total según el precio por libra configurado, completándose en menos de 1 segundo.}{Propietaria -- Bloque 6, Pregunta 29}{Must}{Estable}{Al ingresar peso y precio por libra, el sistema calcula y muestra el costo total en menos de 1 segundo.}{Dependencias: RF-03. EV-04.}

\rfitem{RF-05}{Emitir comprobante de pesaje}{El sistema deberá generar un comprobante con peso, precio por libra y costo total, imprimible o exportable en formato digital.}{Propietaria -- Bloque 6, Pregunta 29}{Should}{Estable}{El comprobante incluye peso, precio unitario, costo total y fecha/hora del registro.}{Dependencias: RF-04. EV-04.}

\rfitem{RF-06}{Registrar despiece y clasificación de cortes}{El sistema deberá permitir registrar los cortes obtenidos al despostar un animal y su peso individual, validando que la suma de pesos no supere el peso del animal más el margen de merma de despiece configurable (valor por defecto: 3\%, rango permitido: 1--8\%).}{Propietaria -- Bloque 5, Preguntas 23 y 25}{Must}{Estable}{La suma de los pesos de los cortes no supera el peso del animal más el margen de merma configurado (defecto 3\%); si supera, el sistema muestra una advertencia y solicita corrección.}{Dependencias: RF-03. EV-05.}

\rfitem{RF-07}{Registrar ingreso de producto a cámara frigorífica}{El sistema deberá permitir registrar el ingreso de un producto a una cámara frigorífica específica, indicando sección (cerdo o res) y fecha de ingreso.}{Propietaria -- Bloque 9, Pregunta 53}{Must}{Estable}{Todo producto almacenado queda asociado a una cámara y una sección válidas.}{Dependencias: RF-06. EV-06.}

\rfitem{RF-08}{Calcular y alertar vida útil del producto}{El sistema deberá calcular los días restantes de vida útil de un producto según su tipo (pollo: 3 días; res: 7 días; cerdo: 5 días; embutidos: 15 días) y generar una alerta visual cuando falte 1 día o menos.}{Propietaria -- Bloque 9, Pregunta 47}{Must}{Volátil}{El sistema muestra una alerta cuando falta 1 día o menos para el límite de vida útil; si el tipo de producto no tiene umbral configurado, genera una notificación de ``sin configurar'' sin interrumpir la operación. La dependencia formal con RF-20 (incidentes de energía) queda documentada para sustentar bajas por cadena de frío.}{Dependencias: RF-07, RF-20 (dependencia formal). EV-06.}

\rfitem{RF-09}{Registrar baja de producto}{El sistema deberá permitir registrar la baja de un producto por vencimiento o mal estado, adjuntando peso, motivo y evidencia fotográfica, con aprobación del administrador general. Una baja no se refleja en el inventario hasta que sea aprobada.}{Propietaria -- Bloque 8, Preguntas 44 y 45}{Must}{Estable}{Una baja de producto no afecta el inventario hasta que el administrador general la aprueba.}{Dependencias: RF-07, RF-20 (dependencia formal añadida en v2.0). EV-07.}

\rfitem{RF-10}{Registrar venta o salida de producto}{El sistema deberá permitir registrar la salida de un producto por venta a un cliente o sucursal, referenciando el número de factura del sistema contable externo.}{Propietaria -- Bloque 10, Pregunta 56}{Must}{Estable}{Toda venta descuenta automáticamente el producto correspondiente del inventario de la cámara de origen.}{Dependencias: RF-07. EV-08.}

\rfitem{RF-11}{Consultar trazabilidad de producto}{El sistema deberá permitir consultar, a partir del número de arete o código de lote, la cadena completa de trazabilidad: granja de origen, guía de origen, faenamiento, pesaje, cámara y venta. Cuando la cadena esté incompleta, el sistema muestra la etiqueta ``En proceso'' en el eslabón faltante.}{Propietaria -- Bloque 3, Pregunta 12}{Must}{Estable}{Ante un número de arete válido, el sistema muestra la cadena completa en menos de 3 segundos; si algún eslabón está incompleto, muestra la etiqueta ``En proceso''.}{Dependencias: RF-02, RF-03, RF-04, RF-06, RF-07, RF-10. EV-03.}

\rfitem{RF-12}{Realizar conteo de inventario}{El sistema deberá permitir registrar un conteo físico semanal de inventario por cámara y comparar contra el inventario del sistema, calculando diferencias en libras y valor monetario.}{Propietaria -- Bloque 8, Pregunta 43}{Should}{Estable}{El sistema muestra automáticamente la diferencia en libras y en valor monetario entre el conteo físico y el inventario registrado.}{Dependencias: RF-07, RF-09. EV-09.}

\rfitem{RF-13}{Generar reportes de ventas, compras y pérdidas}{El sistema deberá generar reportes filtrables por rango de fechas y tipo de producto, mostrando totales de ventas, compras y pérdidas por baja. Los reportes también son filtrables por proveedor (derivado en RF-27).}{Propietaria -- Bloque 10, Pregunta 55}{Must}{Estable}{El reporte generado puede exportarse en PDF y refleja el mismo total que la suma de los registros individuales del período consultado.}{Dependencias: RF-09, RF-10. EV-08.}

\rfitem{RF-14}{Visualizar panel resumen del negocio}{El sistema deberá presentar un panel gráfico con las ventas por tipo de producto, el cliente que más compró y el producto de mayor rotación en un período determinado.}{Propietaria -- Bloque 10, Pregunta 56}{Should}{Volátil}{El panel se actualiza automáticamente al registrarse una nueva venta, sin requerir recarga manual.}{Dependencias: RF-10, RF-13. EV-08.}

\rfitem{RF-15}{Gestionar usuarios y roles}{El sistema deberá permitir crear, editar y desactivar usuarios, asignándoles uno de los roles definidos (propietaria, administrador general, encargado de bodega, carnicero). La autenticación MFA (TOTP) es obligatoria para roles con capacidad de aprobar bajas y configurar precios (RFC-02 aprobado).}{Propietaria -- Bloque 1, Preguntas 2 y 3}{Must}{Estable}{Un usuario desactivado no puede iniciar sesión ni realizar ninguna operación. Los roles con autorización de baja o configuración de precios requieren MFA (TOTP), con pregunta de seguridad como fallback en modo offline (RF-21).}{Dependencias: Ninguna. EV-01.}

\rfitem{RF-16}{Registrar movimientos de inventario por fecha}{El sistema deberá registrar cada ingreso y egreso de producto por cámara y fecha, permitiendo consultar el stock disponible en un momento determinado. Las operaciones offline pendientes de sincronización no se incluyen en el cálculo de stock visible hasta que se sincronicen.}{Propietaria -- Pregunta complementaria de reporte de movimientos}{Must}{Estable}{El sistema permite consultar el stock de un producto a una fecha dada con 0\% de margen de error respecto a los movimientos registrados y sincronizados hasta esa fecha.}{Dependencias: RF-07, RF-09, RF-10. EV-09.}

\rfitem{RF-17}{Registrar georreferencia de proveedor/granja}{El sistema deberá permitir asociar coordenadas geográficas de la granja de origen al perfil del proveedor, como flujo alternativo del CU-01. El campo es opcional pero, si se ingresa, valida rango de latitud/longitud de Ecuador continental con precisión mínima de 10 metros ($\pm$0,0001 grados decimales).}{Segunda ronda -- Administrador general}{Should}{Estable}{Un proveedor sin coordenadas opera normalmente; si se ingresan, el sistema valida el rango geográfico de Ecuador continental y la precisión mínima.}{Dependencias: RF-01. EV-10.}

\rfitem{RF-18}{Consultar historial de precios por libra}{El sistema deberá permitir consultar el histórico de precios por libra configurados para cada tipo de producto en un rango de fechas, mostrando fecha de cambio, precio anterior, precio nuevo y usuario que lo modificó.}{Segunda ronda -- Propietaria}{Could}{Estable}{El histórico muestra, para cada cambio de precio, la fecha, el precio anterior, el precio nuevo y el usuario que lo modificó.}{Dependencias: RF-04. EV-11.}

\rfitem{RF-19}{Configurar precios por tipo de producto y temporada}{El sistema deberá permitir configurar el precio por libra vigente para cada tipo de producto, con posibilidad de definir una fecha de vigencia futura. No se permiten dos precios futuros con fechas de vigencia solapadas para el mismo tipo de producto.}{Segunda ronda -- Propietaria}{Must}{Estable}{Un precio con fecha de vigencia futura no se aplica a pesajes antes de esa fecha; el sistema rechaza la configuración de dos precios futuros solapados para el mismo tipo de producto.}{Dependencias: RF-04. EV-11.}

\rfitem{RF-20}{Registrar incidentes de corte de energía}{El sistema deberá permitir registrar un incidente de corte de energía, indicando fecha, hora de inicio, hora de restablecimiento y cámaras afectadas, como insumo para sustentar posibles bajas de producto.}{Segunda ronda -- Propietaria}{Must}{Estable}{Todo incidente registrado queda disponible como evidencia asociable a una baja de producto (RF-09) generada por pérdida de cadena de frío.}{Dependencias: RF-07. EV-12.}

\rfitem{RF-21}{Operar en modo de contingencia offline}{El sistema deberá permitir registrar operaciones críticas de bodega (ingreso, pesaje) sin conexión a internet y sincronizarlas automáticamente al restablecerse la conectividad, en un máximo de 5 minutos. El modo offline se sostiene al menos 8 horas continuas. El fallback de MFA (pregunta de seguridad) está disponible en modo offline para los roles que lo requieren (RF-15).}{Segunda ronda -- Encargado de bodega}{Must}{Volátil}{Toda operación registrada en modo offline se sincroniza sin pérdida de datos dentro de los 5 minutos posteriores a la reconexión.}{Dependencias: RF-03, RF-04, RF-15. EV-13.}

\rfitem{RF-22}{Registrar devoluciones de clientes}{El sistema deberá permitir registrar la devolución total o parcial de un producto vendido, reingresándolo al inventario o dándolo de baja según su estado. El contador de vida útil del producto devuelto se conserva desde la fecha de ingreso original a cámara.}{Segunda ronda -- Encargado de bodega}{Should}{Estable}{Toda devolución queda vinculada a la venta original mediante el número de factura; el contador de vida útil del producto conserva la fecha de ingreso original.}{Dependencias: RF-10, RF-09. EV-14.}

\rfitem{RF-23}{Exportar evidencia de trazabilidad para auditoría}{El sistema deberá permitir exportar, en PDF o CSV, la cadena de trazabilidad completa de un lote o animal para presentarla ante auditorías de Agrocalidad. El PDF incluye la firma veterinaria digitalizada (RFC-03 aprobado).}{Segunda ronda -- Administrador general}{Must}{Estable}{El archivo exportado contiene los mismos datos mostrados en la consulta de trazabilidad (RF-11), la firma veterinaria digitalizada, y conserva el formato exigido por Agrocalidad.}{Dependencias: RF-11. EV-15.}

\rfitem{RF-24}{Buscar producto por código de lote en cámaras}{El sistema deberá permitir buscar y filtrar productos almacenados en cámaras frigoríficas por código de lote, tipo de producto o sección, retornando resultados en menos de 2 segundos.}{Segunda ronda -- Encargado de bodega}{Should}{Estable}{Una búsqueda por código de lote válido retorna el producto en menos de 2 segundos.}{Dependencias: RF-07. EV-09.}

\rfitem{RF-25}{Generar código QR de trazabilidad}{El sistema deberá generar un código QR asociado a cada lote o animal que, al ser escaneado, dirija a la vista pública de trazabilidad con los datos autorizados para consulta externa (campos delimitados según RF-26).}{Segunda ronda -- Cliente / Propietaria}{Could}{Estable}{El código QR generado enlaza correctamente a la vista pública de trazabilidad del lote/animal correspondiente y muestra únicamente los campos públicos definidos en RF-26.}{Dependencias: RF-11, RF-26. EV-03.}

\rfitem{RF-26}{Vista pública de trazabilidad vía QR -- campos delimitados}{El sistema deberá exponer, a través del enlace generado por el código QR (RF-25), únicamente los siguientes campos de la cadena de trazabilidad: nombre del proveedor, tipo de producto, fecha de ingreso a cámara, fecha de venta y número de guía de origen. No se exponen costos internos, datos personales de clientes ni información de inventario.}{CCB -- RFC implícito D-01 aprobado}{Must}{Estable}{Al escanear el QR, el sistema muestra solo los campos públicos definidos; si se intenta acceder a campos internos por URL directa, el sistema retorna un error 403.}{Dependencias: RF-11, RF-25. CCB/D-01.}

\rfitem{RF-27}{Filtrar reportes por proveedor}{El sistema deberá permitir filtrar los reportes gerenciales (RF-13) por proveedor, generando un desglose de ventas, compras y pérdidas asociadas a cada proveedor dentro del período seleccionado.}{CCB -- RFC-01 aprobado}{Must}{Estable}{Dado un filtro de proveedor activo, cuando se genera el reporte, entonces el sistema muestra únicamente los registros asociados a ese proveedor con el total correcto.}{Dependencias: RF-13, RF-01. CCB/RFC-01.}

%----------------------------------------------------------
\subsection{Requisitos no funcionales}
%----------------------------------------------------------

Los 15 RNF se organizan según las características de calidad de ISO/IEC~25010:2023 \cite{iso25010}.

\begin{longtable}{|p{1.4cm}|p{3cm}|p{9.5cm}|}
\caption{Requisitos no funcionales (RNF-01 a RNF-15) -- ISO/IEC 25010:2023}\\
\hline
\textbf{ID} & \textbf{Característica ISO 25010} & \textbf{Descripción cuantificable} \\
\hline
RNF-01 & Eficiencia de desempeño & El cálculo del costo total al registrar un pesaje deberá completarse en menos de 1 segundo bajo condiciones normales ($\leq$5 usuarios concurrentes). \\
\hline
RNF-02 & Fiabilidad & El sistema deberá estar disponible al menos el 98\% del tiempo mensual. RTO de 30 minutos en horario de faenamiento y 2 horas fuera de ese horario. \\
\hline
RNF-03 & Seguridad & JWT con expiración de 8 horas; MFA mediante TOTP para roles con capacidad de aprobar bajas y configurar precios (RFC-02 aprobado), con pregunta de seguridad como fallback offline. \\
\hline
RNF-04 & Interacción de usuario (Usabilidad) & Un encargado de bodega sin experiencia previa deberá poder completar el registro de un ingreso en un máximo de 5 minutos tras una capacitación de 30 minutos. \\
\hline
RNF-05 & Mantenibilidad & El sistema deberá seguir arquitectura en capas (presentación, lógica de negocio, acceso a datos) que permita agregar un nuevo tipo de reporte sin modificar más del 10\% de los módulos adyacentes (pesaje y trazabilidad). \\
\hline
RNF-06 & Portabilidad & La interfaz deberá ser accesible desde Chrome, Edge, Firefox (dos últimas versiones mayores) sin instalación adicional. \\
\hline
RNF-07 & Seguridad (integridad de datos) & Todo registro de baja de producto queda auditado con usuario, fecha y hora; no podrá eliminarse una vez aprobado. \\
\hline
RNF-08 & Adecuación funcional (completitud) & El conjunto de RF \textit{Debe tener} implementado en el MVP deberá cubrir al menos el 80\% de los RF Must del ERS (mínimo 14 de 17 RF Must originales), verificable mediante la matriz de trazabilidad. \\
\hline
RNF-09 & Compatibilidad & El módulo de pesaje deberá poder integrarse mediante interfaz serial/USB ASCII (protocolo Ohaus) con al menos un modelo de báscula electrónica comercial sin requerir cambios en el modelo de datos. \\
\hline
RNF-10 & Flexibilidad (adaptabilidad) & El sistema deberá permitir agregar un nuevo tipo de producto y su límite de vida útil sin requerir un nuevo despliegue del software, mediante configuración paramétrica. \\
\hline
RNF-11 & Eficiencia de desempeño & La generación de un reporte filtrado por rango de hasta 90 días deberá completarse en menos de 5 segundos con hasta 10\,000 registros. \\
\hline
RNF-12 & Seguridad (confidencialidad) & Datos personales de proveedores, clientes y usuarios almacenados cifrados en reposo (AES-256) y transmitidos únicamente mediante HTTPS; llaves gestionadas mediante KMS con rotación cada 90 días. \\
\hline
RNF-13 & Interacción de usuario (accesibilidad) & Pantallas críticas de bodega con contraste mínimo de 4.5:1 y controles táctiles de al menos 44$\times$44 px. \\
\hline
RNF-14 & Fiabilidad (recuperabilidad) & Ante corte de energía o pérdida de conectividad, el sistema deberá conservar sin pérdida las operaciones en modo offline (RF-21) y sincronizarlas al reiniciar en un máximo de 5 minutos. El modo offline se sostiene al menos 8 horas continuas. \\
\hline
RNF-15 & Mantenibilidad (capacidad de prueba) & El código fuente del MVP deberá contar con cobertura de pruebas automatizadas de al menos el 40\% de las líneas de los módulos de pesaje y trazabilidad. \\
\hline
\end{longtable}

%----------------------------------------------------------
\subsection{Requisito de explicabilidad (componentes de IA)}
%----------------------------------------------------------

En la Entrega 3 (2A) el requisito de explicabilidad se declaró no aplicable por ausencia de componentes IA. En esta Entrega 4 (2B), la PE5 incorpora dos componentes de IA (IA-01 e IA-02) cuyos requisitos de explicabilidad se especifican en la Sección~8. Las dimensiones de explicabilidad aplicadas siguen el marco de \cite{chazette2022explainable}: explicación en lenguaje natural, identificación de los factores más influyentes y tiempo de respuesta de la explicación $\leq 2$ segundos.

%----------------------------------------------------------
\subsection{Requisitos legales}
%----------------------------------------------------------

Trazabilidad explícita entre la Ley Orgánica de Protección de Datos Personales \cite{lopd2021} y los RF/RNF que la operacionalizan.

\begin{longtable}{|p{2.5cm}|p{6cm}|p{5cm}|}
\caption{Trazabilidad de requisitos legales -- LOPDP y RF/RNF}\\
\hline
\textbf{Artículo LOPDP} & \textbf{Obligación} & \textbf{RF/RNF que la satisface} \\
\hline
Art. 7 -- Principio de licitud & Todo tratamiento de datos personales requiere base de licitud (consentimiento o ejecución contractual). & RF-01, RF-15 \\
\hline
Art. 9 -- Deber de información & Las personas titulares deben ser informadas sobre el tratamiento de sus datos al registrarse. & RF-15 \\
\hline
Art. 30 -- Medidas de seguridad & Implementar medidas técnicas y organizativas apropiadas. & RNF-03, RNF-12 \\
\hline
Art. 38 -- Conservación y eliminación & Conservar datos solo el tiempo necesario; eliminarlos de forma segura al finalizar. & RNF-07, RF-13 \\
\hline
Art. 45--49 -- Derechos de titulares & Identificar, corregir o eliminar datos personales a solicitud del titular. & RF-15, RF-01 \\
\hline
Reglamento LOPDP -- Art. 26 & Minimización de datos: RF-26 delimita explícitamente los campos públicos del QR, excluyendo costos internos y datos personales. & RF-26 \\
\hline
HL7 FHIR (aplicabilidad dominio) & No aplica directamente al dominio cárnico; se aplica en el equipo PFC-IR (SIGCM). & -- \\
\hline
\end{longtable}

%----------------------------------------------------------
\subsection{Historias de usuario (formato Connextra) y criterios de aceptación (Gherkin)}
%----------------------------------------------------------

Se redactan historias de usuario para cada RF de prioridad \textit{Must}, verificadas con criterios INVEST y con escenarios Gherkin (Dado--Cuando--Entonces).

\newcommand{\huitem}[4]{%
\noindent\textbf{#1} (trazada a #2)\\
\textit{Como} #3\\[0.15cm]
\textbf{Criterio de aceptación (Gherkin):}\\
#4
\vspace{0.3cm}
}

\huitem{HU-01}{RF-01}{administradora del negocio, quiero registrar un proveedor con su RUC validado, para asegurar que solo trabajo con proveedores regularizados ante el SRI.}{
\textit{Dado} que ingreso al módulo de proveedores,\\
\textit{Cuando} intento guardar un proveedor con un RUC que no supera la verificación de dígito módulo 11,\\
\textit{Entonces} el sistema rechaza el registro y muestra el mensaje de error específico.}

\huitem{HU-02}{RF-02}{encargado de bodega, quiero registrar la guía de origen de un lote con firma veterinaria, para cumplir con la trazabilidad de Agrocalidad antes de aceptar el ingreso.}{
\textit{Dado} que un proveedor no tiene guía de origen con firma veterinaria registrada,\\
\textit{Cuando} intento registrar el ingreso de un lote asociado a ese proveedor,\\
\textit{Entonces} el sistema rechaza el registro del lote hasta que exista una guía con firma veterinaria.}

\huitem{HU-03}{RF-03}{encargado de bodega, quiero registrar el ingreso de animales por número de arete, para mantener la trazabilidad individual exigida por Agrocalidad.}{
\textit{Dado} que registro el ingreso de un cerdo o res,\\
\textit{Cuando} el número de arete ingresado ya existe en el sistema,\\
\textit{Entonces} el sistema rechaza el registro y muestra un mensaje de duplicidad.}

\huitem{HU-04}{RF-04}{encargado de bodega, quiero que el sistema calcule automáticamente el costo total al pesar un producto, para evitar errores de cálculo manual.}{
\textit{Dado} que ingreso el peso en libras y el precio por libra configurado,\\
\textit{Cuando} confirmo el registro de pesaje,\\
\textit{Entonces} el sistema muestra el costo total calculado en menos de 1 segundo.}

\huitem{HU-05}{RF-06}{carnicero, quiero asignar una categoría de venta a cada corte clasificado, para que el producto quede listo para su ingreso a cámara.}{
\textit{Dado} que un corte ya registrado no corresponde a ninguna categoría de venta vigente,\\
\textit{Cuando} intento clasificarlo,\\
\textit{Entonces} el sistema lo marca como ``por definir'' y bloquea su ingreso a cámara hasta que un administrador complete la categoría.}

\huitem{HU-06}{RF-06}{carnicero, quiero registrar el peso de cada corte con margen de merma configurable, para que el inventario refleje la clasificación real de productos.}{
\textit{Dado} que registro los cortes de un animal ya pesado,\\
\textit{Cuando} la suma de los pesos de los cortes supera el peso del animal más el margen de merma configurado (defecto 3\%),\\
\textit{Entonces} el sistema muestra una advertencia y solicita confirmación o corrección.}

\huitem{HU-07}{RF-07}{encargado de bodega, quiero registrar el ingreso de un producto a una cámara específica, para mantener el inventario organizado por sección.}{
\textit{Dado} que un producto pasó por despiece o pesaje,\\
\textit{Cuando} lo registro como ingresado a una cámara sin indicar sección válida,\\
\textit{Entonces} el sistema no permite completar el registro.}

\huitem{HU-08}{RF-08}{administrador general, quiero recibir una alerta cuando un producto esté por vencer con umbral de 1 día, para priorizar su venta o registrar su baja a tiempo.}{
\textit{Dado} que un producto está almacenado en cámara,\\
\textit{Cuando} falta 1 día o menos para el límite de vida útil configurado para su tipo,\\
\textit{Entonces} el sistema muestra una alerta visible en el panel del administrador general.}

\huitem{HU-09}{RF-09}{administrador general, quiero aprobar las bajas de producto antes de que afecten el inventario, para sustentar las pérdidas con evidencia fotográfica.}{
\textit{Dado} que un encargado de bodega registra una baja de producto,\\
\textit{Cuando} la baja aún no ha sido aprobada por el administrador general,\\
\textit{Entonces} el inventario del sistema no refleja la baja.}

\huitem{HU-10}{RF-10}{encargado de bodega, quiero registrar la venta de un producto referenciando el número de factura externo, para mantener consistencia entre el inventario y la facturación.}{
\textit{Dado} que un cliente compra un producto disponible en inventario,\\
\textit{Cuando} confirmo el registro de la venta,\\
\textit{Entonces} el sistema descuenta automáticamente el producto de la cámara de origen.}

\huitem{HU-11}{RF-11}{propietaria, quiero consultar la trazabilidad completa con etiqueta ``En proceso'' cuando la cadena esté incompleta, para responder ante un requerimiento de Agrocalidad.}{
\textit{Dado} que ingreso un número de arete válido,\\
\textit{Cuando} solicito la consulta de trazabilidad,\\
\textit{Entonces} el sistema muestra la cadena completa en menos de 3 segundos; si algún eslabón está incompleto, muestra la etiqueta ``En proceso''.}

\huitem{HU-12}{RF-12}{administrador general, quiero que el sistema me alerte cuando la diferencia del conteo físico supere un umbral configurado, para investigar la desviación antes de cerrar el conteo semanal.}{
\textit{Dado} que registro el conteo físico de una cámara,\\
\textit{Cuando} la diferencia contra el inventario registrado supera el umbral configurado (2\% en libras o en valor monetario),\\
\textit{Entonces} el sistema marca el conteo como ``diferencia significativa'' y exige una justificación antes de cerrarlo.}

\huitem{HU-13}{RF-13}{propietaria, quiero generar reportes de ventas, compras y pérdidas por período, para tomar decisiones informadas de compra y venta.}{
\textit{Dado} que selecciono un rango de fechas y un tipo de producto,\\
\textit{Cuando} genero el reporte,\\
\textit{Entonces} el sistema lo exporta en PDF con un total igual a la suma de los registros individuales del período.}

\huitem{HU-15}{RF-15}{administrador general, quiero gestionar usuarios con MFA obligatoria para roles críticos, para evitar accesos no autorizados y cumplir con la LOPDP.}{
\textit{Dado} que desactivo el usuario de un empleado,\\
\textit{Cuando} ese empleado intenta iniciar sesión,\\
\textit{Entonces} el sistema rechaza el acceso. Y dado un rol con aprobación de bajas, cuando inicia sesión, entonces el sistema solicita MFA (TOTP).}

\huitem{HU-16}{RF-16}{administrador general, quiero consultar el stock a una fecha pasada sin incluir operaciones offline no sincronizadas, para verificar diferencias de inventario reportadas por bodega.}{
\textit{Dado} que existen movimientos registrados y sincronizados para un producto,\\
\textit{Cuando} consulto el stock a una fecha específica,\\
\textit{Entonces} el sistema muestra el stock con 0\% de margen de error respecto a los movimientos sincronizados hasta esa fecha.}

\huitem{HU-19}{RF-19}{propietaria, quiero configurar el precio por libra con vigencia futura sin solapamiento, para planificar cambios de precio por temporada.}{
\textit{Dado} que configuro un nuevo precio con vigencia a partir de una fecha futura,\\
\textit{Cuando} se registra un pesaje antes de esa fecha,\\
\textit{Entonces} el sistema aplica el precio anterior; el sistema rechaza la configuración de dos precios futuros solapados para el mismo tipo de producto.}

\huitem{HU-20}{RF-20}{administrador general, quiero registrar un corte de energía eléctrica para sustentar bajas por pérdida de cadena de frío.}{
\textit{Dado} que registro un incidente de corte de energía con hora de inicio y de restablecimiento,\\
\textit{Cuando} registro posteriormente una baja de producto en una cámara afectada,\\
\textit{Entonces} el sistema permite asociar el incidente como evidencia de la baja.}

\huitem{HU-21}{RF-21}{encargado de bodega, quiero registrar ingresos y pesajes en modo offline con sincronización en máximo 5 minutos, para no detener la operación durante un corte de conectividad.}{
\textit{Dado} que el dispositivo pierde conectividad,\\
\textit{Cuando} registro un ingreso o un pesaje en modo offline,\\
\textit{Entonces} el sistema sincroniza la operación sin pérdida de datos dentro de los 5 minutos posteriores a la reconexión.}

\huitem{HU-23}{RF-23}{administrador general, quiero exportar la trazabilidad completa con firma veterinaria en PDF, para presentarla ante una auditoría de Agrocalidad.}{
\textit{Dado} que un lote tiene trazabilidad completa registrada,\\
\textit{Cuando} exporto la evidencia de trazabilidad,\\
\textit{Entonces} el sistema genera un PDF con los mismos datos de la consulta de trazabilidad, incluyendo la firma veterinaria digitalizada.}

\huitem{HU-26}{RF-26}{cliente final, quiero ver solo los campos públicos de trazabilidad al escanear el QR, para conocer el origen del producto sin exponer datos internos del negocio.}{
\textit{Dado} que escaneo el código QR de un producto,\\
\textit{Cuando} accedo a la vista pública de trazabilidad,\\
\textit{Entonces} el sistema muestra solo: nombre del proveedor, tipo de producto, fecha de ingreso a cámara, fecha de venta y número de guía de origen; el acceso directo a campos internos retorna error 403.}

\huitem{HU-27}{RF-27}{propietaria, quiero filtrar los reportes por proveedor para analizar el desempeño de cada uno.}{
\textit{Dado} que selecciono un proveedor como filtro en el módulo de reportes,\\
\textit{Cuando} genero el reporte,\\
\textit{Entonces} el sistema muestra únicamente los registros asociados a ese proveedor con el total correcto.}

%----------------------------------------------------------
\subsection{Restricciones de diseño}
%----------------------------------------------------------

\begin{enumerate}[nosep]
\item \textbf{RD-1:} Identificación dual: bovinos y porcinos por arete individual; aves por lote. Exigencia de Agrocalidad.
\item \textbf{RD-2:} Sin facturación electrónica SRI (solo referencia al número de factura externo).
\item \textbf{RD-3:} Continuidad offline obligatoria $\geq$8 horas (RF-21, RNF-14).
\item \textbf{RD-4:} Motor de base de datos relacional (PostgreSQL 15+).
\item \textbf{RD-5:} La configuración de precios (RF-19) y de límites de vida útil (RD-2) se almacena en tablas paramétricas editables sin intervención del equipo de desarrollo.
\end{enumerate}

%==========================================================
\section{Modelado del Sistema con UML}
%==========================================================

\subsection{Diagrama general de casos de uso}
\begin{figure}[H]
    \centering
    \includegraphics[width=\textwidth,keepaspectratio]{figures/Use_case_Diagram.png}
    \caption{Diagrama de casos de uso general del sistema CarniCore}
    \label{fig:use_case}
\end{figure}

\subsection{Especificación textual de casos de uso (12 CU detallados)}

\newcommand{\cuitem}[6]{%
\noindent\textbf{#1 -- #2}\\
\textbf{Actor principal:} #3\\
\textbf{Precondiciones:} #4\\
\textbf{Flujo principal:} #5\\
\textbf{Postcondiciones:} #6
\vspace{0.35cm}
}

\cuitem{CU-01}{Registrar ingreso de animal/lote}{Encargado de bodega}{Existe una guía de origen registrada para el proveedor (CU-02).}{(1) El encargado selecciona ``Registrar ingreso''. (2) El sistema solicita proveedor, guía de origen y tipo de producto. (3) Para cerdo/res solicita arete por animal; para aves, la cantidad. (4) El sistema verifica la firma veterinaria de la guía (RF-02) y la georreferencia del proveedor, si existe (RF-17). (5) El encargado confirma. (6) El sistema registra el lote/animal como disponible.}{El lote/animal queda disponible para el pesaje (CU-03).}
\noindent\textbf{\mbox{Flujo alternativo 1}:} el arete no existe en el histórico del proveedor para el animal declarado (cerdo/res, paso 3). (a) El sistema no encuentra el arete registrado. (b) El encargado confirma con el proveedor y registra el arete como nuevo. (c) El sistema da de alta el arete y continúa el flujo principal en el paso 5.\\
\textbf{Poscondición (flujo alternativo 1):} el arete queda dado de alta y disponible para futuros ingresos del mismo proveedor.\\
\noindent\textbf{\mbox{Flujo alternativo 2}:} el proveedor tiene georreferencia registrada (RF-17). (a) El sistema muestra la ubicación del proveedor en el mapa durante el paso 4. (b) El encargado puede usar la ubicación como apoyo para validar la coherencia del origen declarado antes de confirmar.\\
\textbf{Poscondición (flujo alternativo 2):} la georreferencia del proveedor queda visible como apoyo a la decisión, sin bloquear el flujo principal.\\
\noindent\textbf{\mbox{Excepción}:} la guía de origen registrada para el proveedor no tiene firma veterinaria válida (RF-02) $\Rightarrow$ el sistema rechaza el registro y solicita completar la guía antes de continuar.\\
\textbf{Poscondición (excepción):} el ingreso no se registra hasta que la guía cuente con firma veterinaria válida.
\vspace{0.35cm}

\cuitem{CU-02}{Registrar guía de origen}{Encargado de bodega}{El encargado cuenta con los datos del proveedor y la guía física del lote.}{(1) El encargado ingresa número de guía, cantidad de canales y confirma la firma del médico veterinario (digitalizada, RFC-03). (2) El sistema invoca ``Validar sellos veterinarios'' (\textit{include}). (3) El sistema almacena la guía.}{La guía queda disponible para asociarse a un ingreso (CU-01).}
\noindent\textbf{\mbox{Flujo alternativo}:} el proveedor aún no está registrado en el sistema (RF-01). (a) El sistema detecta que el proveedor no existe. (b) El encargado registra el proveedor mínimo (nombre, identificación) antes de continuar. (c) El sistema reanuda el registro de la guía con el proveedor recién creado.\\
\textbf{Poscondición (flujo alternativo):} la guía queda asociada a un proveedor válido.\\
\noindent\textbf{\mbox{Excepción}:} el servicio ``Validar sellos veterinarios'' (\textit{include}) no responde o el sello no coincide con el registro del veterinario $\Rightarrow$ el sistema marca la guía como ``pendiente de validación''.\\
\textbf{Poscondición (excepción):} la guía permanece pendiente hasta revalidación manual del administrador.
\vspace{0.35cm}

\cuitem{CU-03}{Registrar pesaje}{Encargado de bodega}{El lote/animal fue registrado previamente (CU-01).}{(1) El encargado selecciona el lote/animal. (2) Invoca ``Calcular costo por libra'' (\textit{include}) con el precio vigente (RF-19; no se aplican precios futuros pendientes). (3) Calcula costo total en $<$1s. (4) El encargado confirma. (5) El sistema emite el comprobante (RF-05).}{El producto queda disponible para despiece (CU-04) o almacenamiento directo.}
\noindent\textbf{\mbox{Flujo alternativo}:} el pesaje capturado está fuera del rango esperado configurado para el tipo de animal/lote. (a) El sistema marca el pesaje como ``fuera de rango''. (b) Solicita al encargado confirmar un re-pesaje o justificar la desviación. (c) Si se confirma, el sistema continúa con ``Calcular costo por libra'' usando el valor confirmado.\\
\textbf{Poscondición (flujo alternativo):} el pesaje queda marcado con la observación de rango y auditado.\\
\noindent\textbf{\mbox{Excepción}:} el precio vigente por libra no está configurado para el tipo de producto (RF-19) $\Rightarrow$ el sistema bloquea la emisión del comprobante hasta que se configure el precio.\\
\textbf{Poscondición (excepción):} no se emite comprobante mientras no exista precio vigente válido.
\vspace{0.35cm}

\cuitem{CU-04}{Registrar despiece/cortes}{Carnicero}{El animal fue pesado (CU-03).}{(1) El carnicero registra cada corte (tipo, peso). (2) El sistema valida que la suma de pesos no exceda el peso del animal más el margen de merma de despiece (defecto 3\%, configurable 1--8\%). (3) El sistema almacena los cortes.}{Los cortes quedan disponibles para clasificación (CU-05).}
\noindent\textbf{\mbox{Flujo alternativo}:} la suma de pesos de los cortes se acerca al límite superior del margen de merma configurado (1--8\%). (a) El sistema advierte al carnicero que el margen está por agotarse. (b) El carnicero decide continuar registrando cortes o cerrar el despiece.\\
\textbf{Poscondición (flujo alternativo):} el despiece continúa dentro del margen permitido, con la advertencia registrada.\\
\noindent\textbf{\mbox{Excepción}:} la suma de pesos excede el peso del animal más el margen de merma $\Rightarrow$ el sistema rechaza el registro del corte y exige corregir el peso o el margen configurado.\\
\textbf{Poscondición (excepción):} el corte no se almacena hasta que la suma sea válida.
\vspace{0.35cm}

\cuitem{CU-05}{Clasificar producto por parte}{Carnicero}{Existen cortes registrados (CU-04).}{(1) El carnicero asigna una categoría de venta a cada corte. (2) El sistema asocia la categoría al corte.}{El producto clasificado queda disponible para ingreso a cámara (CU-06).}
\noindent\textbf{\mbox{Flujo alternativo}:} el corte no corresponde a ninguna categoría de venta vigente. (a) El sistema no encuentra categoría aplicable. (b) El carnicero asigna la categoría ``por definir'' temporalmente. (c) Un administrador completa la categorización antes del ingreso a cámara (CU-06).\\
\textbf{Poscondición (flujo alternativo):} el corte permanece clasificado como pendiente hasta su categorización definitiva.\\
\noindent\textbf{\mbox{Excepción}:} se intenta asignar dos categorías simultáneas al mismo corte $\Rightarrow$ el sistema rechaza la doble asignación y conserva la última categoría válida.\\
\textbf{Poscondición (excepción):} cada corte conserva una única categoría vigente.
\vspace{0.35cm}

\cuitem{CU-06}{Registrar ingreso a cámara}{Encargado de bodega}{Existe producto clasificado o pesado disponible.}{(1) El encargado selecciona el producto y la cámara/sección destino. (2) El sistema registra el ingreso con fecha y calcula la vida útil inicial según el tipo. (3) El sistema actualiza el inventario.}{El producto queda disponible para control de vida útil (CU-07), venta (CU-09) o conteo (CU-11).}
\noindent\textbf{\mbox{Flujo alternativo}:} la cámara/sección destino está a capacidad máxima configurada. (a) El sistema advierte que no hay espacio disponible. (b) El encargado selecciona una cámara/sección alterna. (c) El sistema recalcula la vida útil inicial si el tipo de cámara lo requiere.\\
\textbf{Poscondición (flujo alternativo):} el producto ingresa a una cámara con capacidad disponible.\\
\noindent\textbf{\mbox{Excepción}:} el tipo de producto no tiene parámetros de vida útil configurados $\Rightarrow$ el sistema registra el ingreso pero marca el producto como ``vida útil sin configurar''.\\
\textbf{Poscondición (excepción):} el producto queda disponible en inventario, pendiente de configuración de vida útil.
\vspace{0.35cm}

\cuitem{CU-07}{Controlar vida útil (frescura)}{Administrador general (con notificación automática del sistema)}{El producto está almacenado en cámara (CU-06).}{(1) El sistema calcula diariamente los días transcurridos desde el ingreso. (2) Si falta 1 día o menos para el límite configurado, genera una alerta en el panel del administrador. (3) El administrador prioriza la venta o registra la baja (CU-10).}{El estado del producto se actualiza a ``próximo a vencer'' o ``vencido''.}
\noindent\textbf{\mbox{Flujo alternativo}:} la caducidad del producto está vencida sin que el administrador haya actuado sobre la alerta previa. (a) El sistema cambia el estado a ``vencido'' y bloquea la venta (CU-09). (b) Notifica al administrador con prioridad alta. (c) El administrador debe iniciar la baja (CU-10) para liberar el bloqueo.\\
\textbf{Poscondición (flujo alternativo):} el producto vencido queda bloqueado para venta hasta su baja formal.\\
\noindent\textbf{\mbox{Excepción}:} el producto no tiene un tipo configurado con parámetros de vida útil $\Rightarrow$ el sistema genera una notificación de ``sin configurar'' en vez de la alerta de vencimiento.\\
\textbf{Poscondición (excepción):} el producto permanece visible en inventario marcado como ``vida útil sin configurar'' hasta que un administrador complete los parámetros del tipo.
\vspace{0.35cm}

\cuitem{CU-08}{Consultar trazabilidad de producto}{Propietaria / Cliente}{El usuario cuenta con un número de arete o código de lote a consultar.}{(1) El usuario ingresa el número de arete o código. (2) El sistema recupera la cadena de trazabilidad. (3) Si hay eslabones incompletos, muestra ``En proceso''. (4) Si el usuario es cliente externo vía QR, muestra solo los campos públicos de RF-26.}{Ninguna (operación de solo lectura).}
\noindent\textbf{\mbox{Flujo alternativo}:} el código de lote corresponde a un producto con trazabilidad compuesta (varios proveedores o cámaras de origen). (a) El sistema recupera todos los eslabones de cada origen. (b) Presenta la cadena consolidada, indicando cada tramo con su proveedor y cámara.\\
\textbf{Poscondición (flujo alternativo):} el usuario visualiza la trazabilidad completa aun cuando el producto proviene de más de un origen.\\
\noindent\textbf{\mbox{Excepción}:} el número de arete o código de lote ingresado no existe en el sistema $\Rightarrow$ el sistema muestra ``código no encontrado'' sin exponer información interna.\\
\textbf{Poscondición (excepción):} no se retorna información de trazabilidad para códigos inválidos.
\vspace{0.35cm}

\cuitem{CU-09}{Registrar salida/venta de producto}{Encargado de bodega}{El producto está disponible en inventario (CU-06).}{(1) El encargado selecciona producto y cantidad. (2) El sistema valida disponibilidad. (3) Solicita el número de factura externo. (4) Descuenta el producto del inventario.}{El inventario y los reportes (RF-13, RF-27) se actualizan.}
\noindent\textbf{\mbox{Flujo alternativo}:} el cliente solicita una cantidad parcial de un lote que ya está próximo a vencer (CU-07). (a) El sistema sugiere priorizar la venta del producto con menor vida útil restante. (b) El encargado puede aceptar la sugerencia o continuar con el producto originalmente seleccionado.\\
\textbf{Poscondición (flujo alternativo):} la venta se registra priorizando, cuando es posible, el producto próximo a vencer.\\
\noindent\textbf{\mbox{Excepción}:} la cantidad solicitada supera el inventario disponible del producto $\Rightarrow$ el sistema rechaza la venta y no descuenta inventario.\\
\textbf{Poscondición (excepción):} el inventario permanece sin cambios y la venta no se registra.
\vspace{0.35cm}

\cuitem{CU-10}{Registrar baja de producto}{Encargado de bodega (solicita) / Administrador general (aprueba)}{El producto está almacenado en cámara; opcionalmente existe un incidente de corte de energía asociado (RF-20).}{(1) El encargado registra peso, motivo y evidencia fotográfica. (2) El sistema envía la solicitud al administrador. (3) El administrador aprueba o rechaza.}{El inventario se actualiza solo tras la aprobación; el registro queda auditado (RNF-07).}
\noindent\textbf{\mbox{Flujo alternativo}:} existe un incidente de corte de energía asociado (RF-20) que afecta al lote. (a) El sistema vincula automáticamente la solicitud de baja al incidente registrado. (b) El administrador revisa el incidente junto con la solicitud antes de aprobar o rechazar.\\
\textbf{Poscondición (flujo alternativo):} la baja queda trazada al incidente de corte de energía que la originó.\\
\noindent\textbf{\mbox{Excepción}:} el administrador rechaza la solicitud de baja $\Rightarrow$ el sistema conserva el producto en inventario y notifica al encargado de bodega el motivo del rechazo.\\
\textbf{Poscondición (excepción):} el inventario no se modifica y la solicitud queda registrada como rechazada, auditada (RNF-07).
\vspace{0.35cm}

\cuitem{CU-11}{Realizar conteo de inventario}{Administrador general}{El sistema tiene al menos una cámara configurada (CU-06).}{(1) El administrador registra el conteo físico semanal por cámara. (2) El sistema compara el conteo contra el inventario registrado. (3) Calcula y muestra la diferencia en libras y en valor monetario.}{Las diferencias quedan documentadas.}
\noindent\textbf{\mbox{Flujo alternativo}:} la diferencia entre el conteo físico y el inventario registrado supera el umbral configurado (2\% en libras o en valor monetario). (a) El sistema marca el conteo como ``diferencia significativa''. (b) Genera una alerta para el administrador general. (c) El administrador debe justificar o investigar la diferencia antes de cerrar el conteo semanal.\\
\textbf{Poscondición (flujo alternativo):} el conteo con diferencia significativa queda pendiente de cierre hasta su justificación.\\
\noindent\textbf{\mbox{Excepción}:} no existe inventario registrado para la cámara seleccionada $\Rightarrow$ el sistema impide iniciar el conteo para esa cámara.\\
\textbf{Poscondición (excepción):} no se genera conteo para cámaras sin inventario previo.
\vspace{0.35cm}

\cuitem{CU-12}{Generar reportes gerenciales}{Propietaria}{El sistema cuenta con al menos un período histórico con ventas, compras o bajas registradas.}{(1) La propietaria selecciona rango de fechas, tipo de producto y opcionalmente proveedor (RF-27). (2) El sistema calcula totales. (3) Exporta el reporte en PDF.}{Ninguna (operación de solo lectura); el reporte queda disponible para descarga.}
\noindent\textbf{\mbox{Flujo alternativo}:} el rango de fechas seleccionado no tiene ventas, compras ni bajas registradas. (a) El sistema informa que no hay datos para el período. (b) Sugiere el rango más cercano con datos disponibles.\\
\textbf{Poscondición (flujo alternativo):} el reporte no se genera; se informa la ausencia de datos.\\
\noindent\textbf{\mbox{Excepción}:} la exportación a PDF falla por error del generador de reportes $\Rightarrow$ el sistema conserva los totales ya calculados y permite reintentar la exportación sin recalcular.\\
\textbf{Poscondición (excepción):} los totales calculados no se pierden ante un fallo de exportación.
\vspace{0.35cm}

\subsection{Diagrama de clases refinado (v2.0)}

El modelo de clases se refina en v2.0 incorporando: \texttt{QrPublico} (RF-26, con campos delimitados), \texttt{FiltroProveedor} (RF-27), \texttt{ComponenteIA} (base para IA-01 e IA-02), \texttt{PredictorDemanda} (IA-01) y \texttt{DetectorAnomalias} (IA-02), además de la gestión KMS para RNF-12.

\begin{figure}[H]
    \centering
    \includegraphics[width=\textwidth,keepaspectratio]{figures/Class-Diagram.png}
    \caption{Diagrama de clases del Sistema CarniCore (versión refinada v2.0)}
    \label{fig:class}
\end{figure}

\subsection{Diagramas de secuencia (CU-01 a CU-12)}

\begin{figure}[H]
    \centering
    \includegraphics[width=\textwidth,keepaspectratio]{figures/01_Sequence_Diagram.png}
    \caption{Diagrama de secuencia -- CU-01: Registrar ingreso de animal/lote}
\end{figure}

\begin{figure}[H]
    \centering
    \includegraphics[width=\textwidth,keepaspectratio]{figures/02_Sequence_Diagram.png}
    \caption{Diagrama de secuencia -- CU-02: Registrar guía de origen}
\end{figure}

\begin{figure}[H]
    \centering
    \includegraphics[width=\textwidth,keepaspectratio]{figures/03_Sequence_Diagram.png}
    \caption{Diagrama de secuencia -- CU-03: Registrar pesaje con cálculo automático de costo}
\end{figure}

\begin{figure}[H]
    \centering
    \includegraphics[width=\textwidth,keepaspectratio]{figures/04_Sequence_Diagram.png}
    \caption{Diagrama de secuencia -- CU-04: Registrar despiece y clasificación}
\end{figure}

\begin{figure}[H]
    \centering
    \includegraphics[width=\textwidth,keepaspectratio]{figures/05_Sequence_Diagram.png}
    \caption{Diagrama de secuencia -- CU-05: Clasificar producto por parte}
\end{figure}

\begin{figure}[H]
    \centering
    \includegraphics[width=\textwidth,keepaspectratio]{figures/06_Sequence_Diagram.png}
    \caption{Diagrama de secuencia -- CU-06: Registrar ingreso a cámara frigorífica}
\end{figure}

\begin{figure}[H]
    \centering
    \includegraphics[width=\textwidth,keepaspectratio]{figures/07_Sequence_Diagram.png}
    \caption{Diagrama de secuencia -- CU-07: Controlar vida útil}
\end{figure}

\begin{figure}[H]
    \centering
    \includegraphics[width=\textwidth,keepaspectratio]{figures/08_Sequence_Diagram.png}
    \caption{Diagrama de secuencia -- CU-08: Consultar trazabilidad de producto}
\end{figure}

\begin{figure}[H]
    \centering
    \includegraphics[width=\textwidth,keepaspectratio]{figures/09_Sequence_Diagram.png}
    \caption{Diagrama de secuencia -- CU-09: Registrar salida/venta de producto}
\end{figure}

\begin{figure}[H]
    \centering
    \includegraphics[width=\textwidth,keepaspectratio]{figures/10_Sequence_Diagram.png}
    \caption{Diagrama de secuencia -- CU-10: Registrar baja de producto con aprobación}
\end{figure}

\begin{figure}[H]
    \centering
    \includegraphics[width=\textwidth,keepaspectratio]{figures/11_Sequence_Diagram.png}
    \caption{Diagrama de secuencia -- CU-11: Realizar conteo de inventario}
\end{figure}

\begin{figure}[H]
    \centering
    \includegraphics[width=\textwidth,keepaspectratio]{figures/12_Sequence_Diagram.png}
    \caption{Diagrama de secuencia -- CU-12: Generar reportes gerenciales}
\end{figure}

\subsection{Diagramas de actividad (flujos principales)}

\begin{figure}[H]
    \centering
    \includegraphics[width=\textwidth,keepaspectratio]{figures/01_Activity_Diagram_CU-01.png}
    \caption{Diagrama de actividad -- CU-01: Registrar ingreso de animal/lote}
\end{figure}

\begin{figure}[H]
    \centering
    \includegraphics[width=\textwidth,keepaspectratio]{figures/02_Activity_Diagram_CU-02.png}
    \caption{Diagrama de actividad -- CU-02: Registrar guía de origen}
\end{figure}

\begin{figure}[H]
    \centering
    \includegraphics[width=\textwidth,keepaspectratio]{figures/03_Activity_Diagram_CU-03.png}
    \caption{Diagrama de actividad -- CU-03: Registrar pesaje}
\end{figure}

\begin{figure}[H]
    \centering
    \includegraphics[width=\textwidth,keepaspectratio]{figures/04_Activity_Diagram_CU-04.png}
    \caption{Diagrama de actividad -- CU-04: Registrar despiece y clasificación}
\end{figure}

\begin{figure}[H]
    \centering
    \includegraphics[width=\textwidth,keepaspectratio]{figures/05_Activity_Diagram_CU-05.png}
    \caption{Diagrama de actividad -- CU-05: Clasificar producto por parte}
\end{figure}

\begin{figure}[H]
    \centering
    \includegraphics[width=\textwidth,keepaspectratio]{figures/06_Activity_Diagram_CU-06.png}
    \caption{Diagrama de actividad -- CU-06: Registrar ingreso a cámara frigorífica}
\end{figure}

\subsection{Diagrama de estados}

La entidad con ciclo de vida no trivial es \texttt{ProductoAlmacenado}: \textit{Disponible} $\rightarrow$ \textit{Próximo a vencer} $\rightarrow$ \{\textit{Vendido} $\mid$ \textit{Dado de baja (pendiente)} $\rightarrow$ \textit{Dado de baja (aprobado)}\} $\rightarrow$ \textit{Devuelto} (retorna a \textit{Disponible} o a baja según estado).

\begin{figure}[H]
    \centering
    \includegraphics[width=0.75\textwidth,keepaspectratio]{figures/State_Diagram.png}
    \caption{Diagrama de estados de \texttt{ProductoAlmacenado}}
    \label{fig:state}
\end{figure}

\subsection{Diagrama de componentes}

\begin{figure}[H]
    \centering
    \includegraphics[width=\textwidth,keepaspectratio]{figures/Component_Diagram (2).png}
    \caption{Diagrama de componentes del sistema CarniCore}
    \label{fig:components}
\end{figure}

\subsection{Diagrama de despliegue}

\begin{figure}[H]
    \centering    \includegraphics[width=0.9\textwidth,keepaspectratio]{figures/Deployement_Diagram (2).png}
    \caption{Diagrama de despliegue del sistema CarniCore}
    \label{fig:deployment}
\end{figure}

%==========================================================
\section{Priorización y Trazabilidad Extendida}
%==========================================================

\subsection{Priorización combinada MoSCoW + Kano + WSJF}

\begin{longtable}{|p{1.3cm}|p{1.8cm}|p{2.2cm}|p{1.6cm}|p{7.5cm}|}
\caption{Priorización combinada MoSCoW + Kano + WSJF de los 27 RF}\\
\hline
\textbf{ID-RF} & \textbf{MoSCoW} & \textbf{Kano} & \textbf{WSJF*} & \textbf{Justificación breve} \\
\hline
RF-01 & Must & Básica & 8.0 & Sin proveedor regularizado no hay trazabilidad legal de origen. \\
RF-02 & Must & Básica & 8.0 & Exigencia directa de Agrocalidad; riesgo regulatorio si se omite. \\
RF-03 & Must & Básica & 8.0 & Núcleo operativo diario de bodega. \\
RF-04 & Must & Básica & 7.5 & Reemplaza el cálculo manual, fuente principal de errores actuales. \\
RF-05 & Should & Desempeño & 3.0 & Mejora la experiencia pero no bloquea la operación. \\
RF-06 & Must & Básica & 7.0 & Sin despiece registrado no hay inventario por corte. \\
RF-07 & Must & Básica & 7.0 & Base del control de inventario por cámara. \\
RF-08 & Must & Deleite & 6.5 & Ataca directamente la pérdida por vencimiento. \\
RF-09 & Must & Básica & 6.5 & Sustenta pérdidas con evidencia y control gerencial. \\
RF-10 & Must & Básica & 7.0 & Sin descuento automático el inventario queda desincronizado. \\
RF-11 & Must & Deleite & 6.0 & Responde a auditorías de Agrocalidad en tiempo real. \\
RF-12 & Should & Desempeño & 2.5 & Mejora el control pero no es indispensable para operar. \\
RF-13 & Must & Desempeño & 5.5 & Insumo directo de decisiones gerenciales. \\
RF-14 & Should & Deleite & 2.0 & Valor agregado gerencial, no bloqueante. \\
RF-15 & Must & Básica & 6.0 & Seguridad y control de acceso mínimo viable. \\
RF-16 & Must & Básica & 5.5 & Permite auditar diferencias de inventario por fecha. \\
RF-17 & Should & Deleite & 1.5 & Enriquece la trazabilidad, no crítico. \\
RF-18 & Could & Indiferente & 1.0 & Consulta histórica de valor secundario. \\
RF-19 & Must & Desempeño & 5.0 & Necesario para operar cambios de precio sin intervención técnica. \\
RF-20 & Must & Deleite & 4.5 & Sustenta bajas por apagones identificados en la entrevista. \\
RF-21 & Must & Básica & 6.0 & Continuidad operativa ante conectividad intermitente. \\
RF-22 & Should & Desempeño & 2.5 & Mejora el control de postventa. \\
RF-23 & Must & Deleite & 4.0 & Reduce tiempo de respuesta ante auditorías. \\
RF-24 & Should & Desempeño & 2.0 & Mejora la usabilidad operativa de bodega. \\
RF-25 & Could & Deleite & 1.0 & Diferenciador de cara al cliente, no bloqueante. \\
RF-26 & Must & Básica & 5.5 & Delimita campos públicos del QR; exigido por LOPDP (minimización). \\
RF-27 & Must & Desempeño & 4.5 & Permite análisis por proveedor; aprobado por CCB (RFC-01). \\
\hline
\end{longtable}
\footnotesize{*WSJF estimado en escala relativa 0--10 (costo de retraso / tamaño estimado).}

\subsection{Matriz de trazabilidad final end-to-end (84 filas: 42 base + 18 IA + 24 flujos)}

La cadena mínima implementada es: Ley/Art. $\to$ Objetivo $\to$ Interesado $\to$ EV $\to$ RF/RNF $\to$ CU $\to$ HU $\to$ CA $\to$ Componente $\to$ Mockup.

La matriz completa en formato CSV se ubica en \texttt{04\_Trazabilidad/Matriz\_Trazabilidad.csv} del repositorio.

\scriptsize
\setlength{\tabcolsep}{2.2pt}
\begin{longtable}{|p{0.3cm}|p{1.0cm}|p{1.9cm}|p{1.3cm}|p{0.6cm}|p{0.7cm}|p{0.6cm}|p{1.1cm}|p{0.6cm}|p{0.6cm}|p{2.0cm}|p{0.8cm}|}
\caption{Matriz de trazabilidad end-to-end (84 filas: 42 base + 18 IA + 24 flujos)}\\
\hline
\textbf{ID} & \textbf{Ley-Art.} & \textbf{Objetivo} & \textbf{Interesado} & \textbf{EV} & \textbf{RF / RNF} & \textbf{CU} & \textbf{Flujo} & \textbf{HU} & \textbf{CA} & \textbf{Componente} & \textbf{Mockup} \\
\hline
1 & -- & Trazabilidad de origen & Propietaria & EV-01 & RF-01 & CU-01 & Ppal. & HU-01 & CA-01 & Backend-Proveedores & MU-04 \\
2 & Art.30 & Certificación sanitaria & Propietaria & EV-02 & RF-02 & CU-02 & Ppal. & HU-02 & CA-02 & Backend-Guias & MU-05 \\
3 & -- & Registro de ingreso & Bodega & EV-03 & RF-03 & CU-01 & Ppal. & HU-03 & CA-03 & Backend-Ingresos & MU-06 \\
4 & -- & Cálculo de costo & Bodega & EV-04 & RF-04 & CU-03 & Ppal. & HU-04 & CA-04 & Backend-Pesaje & MU-07 \\
5 & -- & Comprobante físico & Bodega & EV-04 & RF-05 & CU-03 & Ppal. & -- & -- & Backend-Pesaje & MU-19 \\
6 & -- & Registro de cortes & Carniceros & EV-05 & RF-06 & CU-04 & Ppal. & HU-06 & CA-06 & Backend-Despiece & MU-08 \\
7 & -- & Control por cámara & Admin. gral. & EV-06 & RF-07 & CU-06 & Ppal. & HU-07 & CA-07 & Backend-Inventario & MU-09 \\
8 & -- & Evitar pérdidas & Admin. gral. & EV-06 & RF-08 & CU-07 & Ppal. & HU-08 & CA-08 & Backend-VidaUtil & MU-11 \\
9 & Art.38 & Sustentar bajas & Admin. gral. & EV-07 & RF-09 & CU-10 & Ppal. & HU-09 & CA-09 & Backend-Bajas & MU-13 \\
10 & -- & Registrar ventas & Bodega / Cliente & EV-08 & RF-10 & CU-09 & Ppal. & HU-10 & CA-10 & Backend-Ventas & MU-14 \\
11 & -- & Historial completo & Agrocalidad & EV-03 & RF-11 & CU-08 & Ppal. & HU-11 & CA-11 & Backend-Trazabilidad & MU-16 \\
12 & -- & Diferencias inventario & Admin. gral. & EV-09 & RF-12 & CU-11 & Ppal. & -- & -- & Backend-Inventario & MU-17 \\
13 & -- & Decisiones gerenciales & Propietaria & EV-08 & RF-13 & CU-12 & Ppal. & HU-13 & CA-13 & Backend-Reportes & MU-18 \\
14 & -- & Desempeño del negocio & Propietaria & EV-08 & RF-14 & CU-12 & Ppal. & -- & -- & Backend-Reportes & MU-02 \\
15 & Art.7,9 & Delegar por rol & Propietaria & EV-01 & RF-15 & -- & -- & HU-15 & CA-15 & Backend-Usuarios & MU-20 \\
16 & -- & Stock por fecha & Propietaria & EV-09 & RF-16 & CU-06 & Ppal. & HU-16 & CA-16 & Backend-Inventario & MU-15 \\
17 & -- & Georreferencia & Admin. gral. & EV-10 & RF-17 & CU-01 & Alt. & -- & -- & Backend-Proveedores & MU-04 \\
18 & -- & Historial de precios & Propietaria & EV-11 & RF-18 & -- & -- & -- & -- & Backend-Precios & MU-14 \\
19 & -- & Configurar precios & Propietaria & EV-11 & RF-19 & CU-03 & Ppal. & HU-19 & CA-19 & Backend-Precios & MU-07 \\
20 & -- & Registrar apagones & Propietaria & EV-12 & RF-20 & CU-10 & Ppal. & HU-20 & CA-20 & Backend-Incidentes & -- \\
21 & -- & Continuidad operativa & Bodega & EV-13 & RF-21 & CU-01/03 & Ppal. & HU-21 & CA-21 & Cliente-Offline & -- \\
22 & -- & Devoluciones & Bodega & EV-14 & RF-22 & CU-09 & Ppal. & -- & -- & Backend-Ventas & MU-14 \\
23 & -- & Auditoría Agrocalidad & Admin. gral. & EV-15 & RF-23 & CU-08 & Ppal. & HU-23 & CA-23 & Backend-Trazabilidad & MU-16 \\
24 & -- & Búsqueda de producto & Bodega & EV-09 & RF-24 & CU-06 & Ppal. & -- & -- & Backend-Inventario & MU-10 \\
25 & -- & QR de trazabilidad & Cliente & EV-03 & RF-25 & CU-08 & Ppal. & -- & -- & Backend-Trazabilidad & MU-16 \\
26 & Art.R26 LOPDP & Campos públicos QR & Cliente & CCB-D01 & RF-26 & CU-08 & Ppal. & HU-26 & CA-26 & Backend-QrPublico & MU-16 \\
27 & -- & Filtro por proveedor & Propietaria & CCB-RFC01 & RF-27 & CU-12 & Ppal. & HU-27 & CA-27 & Backend-Reportes & MU-18 \\
28 & -- & Rendimiento pesaje & Bodega & EV-04 & RNF-01 & CU-03 & Ppal. & -- & -- & Backend-Pesaje & -- \\
29 & -- & Disponibilidad & Todos & -- & RNF-02 & -- & -- & -- & -- & Infraestructura & -- \\
30 & Art.30 & Autenticación segura & Admin. gral. & EV-01 & RNF-03 & -- & -- & -- & -- & Backend-Usuarios & -- \\
31 & -- & Usabilidad bodega & Bodega & EV-09 & RNF-04 & CU-01 & Ppal. & -- & -- & Cliente-Web/Móvil & -- \\
32 & -- & Arquitectura mantenible & Equipo & -- & RNF-05 & -- & -- & -- & -- & Backend (todos) & -- \\
33 & -- & Acceso multi-navegador & Todos & -- & RNF-06 & -- & -- & -- & -- & Cliente-Web & -- \\
34 & Art.38 & Auditoría de bajas & Admin. gral. & EV-07 & RNF-07 & CU-10 & Ppal. & -- & -- & Backend-Bajas & -- \\
35 & -- & Cobertura RF Must MVP & Docente & EV-16 & RNF-08 & -- & -- & -- & -- & MVP & -- \\
36 & -- & Integración báscula & Bodega & -- & RNF-09 & CU-03 & Ppal. & -- & -- & Backend-Pesaje & -- \\
37 & -- & Nuevos tipos producto & Propietaria & EV-11 & RNF-10 & -- & -- & -- & -- & Backend-Precios & -- \\
38 & -- & Reportes alto volumen & Propietaria & EV-08 & RNF-11 & CU-12 & Ppal. & -- & -- & Backend-Reportes & -- \\
39 & Art.30 & Cifrado datos & Todos & -- & RNF-12 & -- & -- & -- & -- & Backend/BD & -- \\
40 & -- & Accesibilidad táctil & Bodega & EV-09 & RNF-13 & CU-01 & Ppal. & -- & -- & Cliente-Móvil & -- \\
41 & -- & Recuperabilidad offline & Bodega & EV-13 & RNF-14 & CU-01/03 & Ppal. & HU-21 & CA-21 & Cliente-Offline & -- \\
42 & -- & Calidad código MVP & Docente & EV-16 & RNF-15 & -- & -- & -- & -- & MVP & -- \\
\hline
\multicolumn{11}{|l|}{\textbf{Extensión IA -- filas 43--60}} \\
\hline
43 & -- & Predicción demanda semanal & Propietaria & EV-08 & RF-IA-01 & -- & -- & -- & -- & IA-PredictorDemanda & -- \\
44 & -- & Intervalo de confianza pred. & Propietaria & EV-08 & RF-IA-02 & -- & -- & -- & -- & IA-PredictorDemanda & -- \\
45 & -- & Fallback predictor & Propietaria & EV-08,12 & RF-IA-03 & -- & -- & -- & -- & IA-PredictorDemanda & -- \\
46 & -- & MAE predictor $\leq$10 lb & Propietaria & EV-08 & RNF-IA-01 & -- & -- & -- & -- & IA-PredictorDemanda & -- \\
47 & -- & Inferencia $<$5s predictor & Propietaria & EV-08 & RNF-IA-02 & -- & -- & -- & -- & IA-PredictorDemanda & -- \\
48 & -- & Escalabilidad predictor & Propietaria & EV-08 & RNF-IA-03 & -- & -- & -- & -- & IA-PredictorDemanda & -- \\
49 & -- & Equidad por temporada & Propietaria & EV-08 & RNF-IA-04 & -- & -- & -- & -- & IA-PredictorDemanda & -- \\
50 & -- & Equidad por tipo corte & Propietaria & EV-08 & RNF-IA-05 & -- & -- & -- & -- & IA-PredictorDemanda & -- \\
51 & -- & Explicabilidad pred. & Propietaria & EV-08 & RNF-IA-06 & -- & -- & -- & -- & IA-PredictorDemanda & -- \\
52 & -- & Clasificar riesgo incidente & Admin. gral. & EV-12 & RF-IA-04 & CU-10 & Ppal. & -- & -- & IA-DetectorAnomalias & -- \\
53 & -- & Prob. y recomendación & Admin. gral. & EV-12 & RF-IA-05 & CU-10 & Ppal. & -- & -- & IA-DetectorAnomalias & -- \\
54 & Art.38 & Registrar decisión admin. & Admin. gral. & EV-12 & RF-IA-06 & CU-10 & Ppal. & -- & -- & IA-DetectorAnomalias & -- \\
55 & -- & Recall $\geq$0.90 detector & Admin. gral. & EV-12 & RNF-IA-07 & -- & -- & -- & -- & IA-DetectorAnomalias & -- \\
56 & -- & Inferencia $<$2s detector & Admin. gral. & EV-12 & RNF-IA-08 & -- & -- & -- & -- & IA-DetectorAnomalias & -- \\
57 & -- & Fallback $<$500ms & Admin. gral. & EV-12 & RNF-IA-09 & -- & -- & -- & -- & IA-DetectorAnomalias & -- \\
58 & -- & Equidad por producto & Admin. gral. & EV-12 & RNF-IA-10 & -- & -- & -- & -- & IA-DetectorAnomalias & -- \\
59 & -- & Equidad por horario & Admin. gral. & EV-12 & RNF-IA-11 & -- & -- & -- & -- & IA-DetectorAnomalias & -- \\
60 & Art.38 & Explicabilidad detector & Admin. gral. & EV-12 & RNF-IA-12 & -- & -- & -- & -- & IA-DetectorAnomalias & -- \\
\hline
\multicolumn{12}{|l|}{\textbf{Extensión flujos alternativos y excepciones -- filas 61--84 (examen suspenso, 16/09/2026)}} \\
\hline
61 &  & Arete inexistente en histórico del proveedor & Bodega & EV-03 & RF-03 & CU-01 & Alt. & HU-03 & CA-03 & Backend-Ingresos & MU-06 \\
62 &  & Guía sin firma veterinaria registrada & Bodega & EV-02 & RF-02 & CU-01 & Exc. & HU-03 & CA-03 & Backend-Guias & MU-05 \\
63 &  & Proveedor no registrado al ingresar guía & Propietaria & EV-01 & RF-01 & CU-02 & Alt. & HU-02 & CA-02 & Backend-Proveedores & MU-04 \\
64 &  & Sello veterinario no válido o servicio no responde & Propietaria & EV-02 & RF-02 & CU-02 & Exc. & HU-02 & CA-02 & Backend-Guias & MU-05 \\
65 &  & Pesaje fuera del rango esperado & Bodega & EV-04 & RF-04 & CU-03 & Alt. & HU-04 & CA-04 & Backend-Pesaje & MU-07 \\
66 &  & Precio vigente no configurado & Propietaria & EV-11 & RF-19 & CU-03 & Exc. & HU-19 & CA-19 & Backend-Precios & MU-07 \\
67 &  & Margen de merma por agotarse & Carniceros & EV-05 & RF-06 & CU-04 & Alt. & HU-06 & CA-06 & Backend-Despiece & MU-08 \\
68 &  & Suma de cortes excede margen de merma & Carniceros & EV-05 & RF-06 & CU-04 & Exc. & HU-06 & CA-06 & Backend-Despiece & MU-08 \\
69 &  & Corte sin categoría de venta vigente & Carniceros & EV-05 & RF-06 & CU-05 & Alt. & HU-05 & CA-05 & Backend-Despiece & MU-08 \\
70 &  & Doble categoría asignada al mismo corte & Carniceros & EV-05 & RF-06 & CU-05 & Exc. & HU-05 & CA-05 & Backend-Despiece & MU-08 \\
71 &  & Cámara/sección destino a capacidad máxima & Admin. gral. & EV-06 & RF-07 & CU-06 & Alt. & HU-07 & CA-07 & Backend-Inventario & MU-09 \\
72 &  & Producto sin parámetros de vida útil configurados & Admin. gral. & EV-06 & RF-08 & CU-06 & Exc. & HU-07 & CA-07 & Backend-VidaUtil & MU-11 \\
73 &  & Caducidad vencida sin acción del administrador & Admin. gral. & EV-06 & RF-08 & CU-07 & Alt. & HU-08 & CA-08 & Backend-VidaUtil & MU-11 \\
74 &  & Producto sin tipo configurado para vida útil & Admin. gral. & EV-06 & RF-08 & CU-07 & Exc. & HU-08 & CA-08 & Backend-VidaUtil & MU-11 \\
75 &  & Trazabilidad compuesta (múltiples orígenes) & Agrocalidad & EV-03 & RF-11 & CU-08 & Alt. & HU-11 & CA-11 & Backend-Trazabilidad & MU-16 \\
76 &  & Código de lote/arete inexistente & Agrocalidad & EV-03 & RF-11 & CU-08 & Exc. & HU-11 & CA-11 & Backend-Trazabilidad & MU-16 \\
77 &  & Priorizar producto próximo a vencer en venta & Bodega & EV-08 & RF-10 & CU-09 & Alt. & HU-10 & CA-10 & Backend-Ventas & MU-14 \\
78 &  & Inventario insuficiente para la venta & Bodega & EV-08 & RF-10 & CU-09 & Exc. & HU-10 & CA-10 & Backend-Ventas & MU-14 \\
79 &  & Baja vinculada a incidente de corte de energía & Propietaria & EV-12 & RF-20 & CU-10 & Alt. & HU-20 & CA-20 & Backend-Incidentes & -- \\
80 &  & Baja rechazada por el administrador & Admin. gral. & EV-07 & RF-09 & CU-10 & Exc. & HU-09 & CA-09 & Backend-Bajas & MU-13 \\
81 &  & Diferencia significativa en conteo de inventario & Admin. gral. & EV-09 & RF-12 & CU-11 & Alt. & HU-12 & CA-12 & Backend-Inventario & MU-17 \\
82 &  & Cámara sin inventario previo registrado & Admin. gral. & EV-09 & RF-12 & CU-11 & Exc. & HU-12 & CA-12 & Backend-Inventario & MU-17 \\
83 &  & Rango de fechas sin datos para reporte & Propietaria & EV-08 & RF-13 & CU-12 & Alt. & HU-13 & CA-13 & Backend-Reportes & MU-18 \\
84 &  & Falla en exportación de reporte a PDF & Propietaria & EV-08 & RF-13 & CU-12 & Exc. & HU-13 & CA-13 & Backend-Reportes & MU-18 \\
\hline
\end{longtable}
\setlength{\tabcolsep}{6pt}
\normalsize

%==========================================================
\section{Producto Mínimo Viable (MVP)}
%==========================================================

El MVP de CarniCore se aloja en el repositorio en la carpeta \texttt{05\_MVP/Ejecutable/Carnicore}, con un \texttt{README.md} de instrucciones para levantar el sistema con \texttt{docker compose up} y un \texttt{video\_demo.mp4} tutorial. El MVP cubre al menos el 80\% de los RF Must (mínimo 14 de los 17 RF Must originales, RF-01 a RF-23 con prioridad Must): RF-01, 02, 03, 04, 06, 07, 08, 09, 10, 11, 13, 15, 16, 19, 20, 21, 23, además de RF-26 y RF-27 aprobados por el CCB.

\begin{figure}[H]
    \centering    \includegraphics[width=0.9\textwidth,keepaspectratio]{figures/MVP_ejecucion.jpeg}
    \caption{MVP CarniCore -- captura de pantalla en ejecución con \texttt{docker compose up}}
    \label{fig:mvp_ejecucion}
\end{figure}

%==========================================================
\section{Componente Empírico -- Diseño del Estudio}
%==========================================================

\subsection{Pregunta de investigación (PICOC)}

El equipo trabaja con el \textbf{Enfoque 2 -- Detección automática de ambigüedad y malos olores en los requisitos}, aplicado sobre el conjunto de 27 RF del ERS/SRS v2.0.

\begin{itemize}[nosep]
\item \textbf{Población:} los 27 RF funcionales de CarniCore (RF-01 a RF-27).
\item \textbf{Intervención:} un detector automático de ambigüedad y malos olores (cuantificadores vagos, conjunciones múltiples, voz pasiva sin agente explícito), implementado en Python con expresiones regulares.
\item \textbf{Comparación:} juicio consensuado de al menos tres personas expertas en IR, clasificando cada RF como ``ambiguo'' o ``no ambiguo'' de forma ciega e independiente.
\item \textbf{Resultado:} precisión, \textit{recall} y medida F1 del detector respecto al consenso experto; $\kappa$ de Cohen \cite{cohen1960kappa} y $\kappa$ de Fleiss \cite{fleiss1971kappa}.
\item \textbf{Contexto:} proyecto académico de IR aplicado a dominio real de distribución cárnica en Ecuador.
\end{itemize}

\textbf{RQ1:} ¿Con qué precisión un detector automático replica el juicio de personas expertas al identificar RF ambiguos en el conjunto de requisitos de CarniCore?

\subsection{Protocolo experimental}

El protocolo completo se encuentra en \texttt{06\_Experimento/protocolo.pdf} (v1.0, 02/08/2026) y fue registrado en OSF antes de toda recolección de datos (\url{https://osf.io/yp7t3}). Las desviaciones del protocolo respecto al pre-registro se documentan en \texttt{06\_Experimento/osf\_deviations.pdf}.

\textbf{Diseño.} Estudio observacional de exactitud diagnóstica, transversal y censal: se aplican dos procedimientos de clasificación (detector automático; consenso experto) sobre la población completa de 27 RF.

\textbf{Consideraciones éticas.} El estudio opera bajo categoría de riesgo C (riesgo mínimo operativo). No se recolectan datos personales identificables más allá de un código de participante. La documentación ética completa reside en \texttt{08\_Etica/} y contiene el paquete base A.1--A.13, el aval institucional y la Adenda de Segunda Ronda.

\textbf{Plan de análisis.} Estadística descriptiva; matriz de confusión; cálculo de precisión, \textit{recall} y F1 con intervalos de confianza; $\kappa$ de Cohen y $\kappa$ de Fleiss. Análisis implementado en Python (\texttt{scikit-learn}, \texttt{statsmodels}) como pipeline reproducible en \texttt{06\_Experimento/scripts\_analisis/} (\texttt{01\_importar\_datos.py}, \texttt{02\_calcular\_kappa.py}, \texttt{03\_matriz\_confusion\_prf1.py}, \texttt{04\_bootstrap\_ic95.py}, orquestados por \texttt{run\_all.py}), con versiones de librerías fijadas en \texttt{requirements.txt}.

\textbf{Resultados de RQ1 (comparación completa detector vs. panel experto).} El detector automático ejecutado sobre los 27 RF no activó ninguna categoría de ambigüedad (0/27, 0\%), resultado consistente con el obtenido en la Entrega 3 sobre los 25 RF originales (también 0/25). El panel de tres personas expertas en IR evaluó, de forma ciega e independiente, los mismos 27 RF y alcanzó consenso (mayoría $\geq$2/3) en marcar 4 de ellos como ambiguos: RF-08, RF-17, RF-21 y RF-22. Al no predecir el detector ningún caso positivo, precisión, \textit{recall} y F1 resultan en 0,0000 (IC 95\% por \textit{bootstrap} no paramétrico, 10\,000 réplicas, semilla fija=42: $[0{,}0,\ 0{,}0]$ en los tres casos). La matriz de confusión es VN=23, FP=0, FN=4, VP=0.

El acuerdo inter-experto fue $\kappa$ de Fleiss $=0{,}2636$ (consenso de 3), y $\kappa$ de Cohen por pareja de $0{,}3478$, $0{,}3571$ y $0{,}0870$ respectivamente -- acuerdo aceptable a leve según Landis y Koch (1977) \cite{cohen1960kappa,fleiss1971kappa}. Este resultado no invalida el estudio: constituye un hallazgo negativo válido para RQ1. El enfoque léxico del detector (cuantificadores vagos, conjunciones múltiples, voz pasiva sin agente) no captura el tipo de ambigüedad que el panel experto identifica en estos 4 RF, de naturaleza más semántica que sintáctica (por ejemplo, RF-08 fue marcado como ambiguo por unanimidad del panel sin activar ninguna de las tres categorías léxicas del detector). El propio acuerdo inter-experto moderado-bajo sugiere además que el techo de acuerdo humano sobre qué constituye ambigüedad limita cuánto puede acertar cualquier sistema automatizado basado en reglas. Dado que N=27 corresponde a la población censal de RF de la v2.0 y no a una muestra de un universo mayor, el intervalo de confianza por \textit{bootstrap} cuantifica la sensibilidad de las métricas a la composición específica de este conjunto, no una generalización estadística externa. El detalle completo de la comparación (matrices de confusión, scripts y datos anonimizados) se encuentra en \texttt{06\_Experimento/resultados/} y \texttt{06\_Experimento/scripts\_analisis/}, reproducible mediante \texttt{python run\_all.py}. Las evaluaciones anonimizadas del panel residen en \texttt{06\_Experimento/datos\_crudos/}; los nombres reales de los evaluadores se conservan exclusivamente en la zona [R] restringida del repositorio, conforme al consentimiento informado firmado.

%==========================================================
\section{Requisitos de Inteligencia Artificial}
%==========================================================

\subsection{Justificación de los componentes de IA}

Si bien el ERS v3.0 declaró que CarniCore no incorporaba componentes de IA en su alcance inicial, la Entrega 4 (2B) exige la especificación de al menos dos componentes plausibles para el dominio. Se identificaron dos con valor demostrable:

\begin{itemize}[nosep]
\item \textbf{IA-01 -- Predictor de Demanda por Tipo de Corte:} anticipa la demanda semanal por tipo de corte, entrenándose con el historial de RF-10 y RF-13.
\item \textbf{IA-02 -- Detector de Anomalías de Cadena de Frío:} clasifica automáticamente el riesgo de pérdida de cadena de frío para cada cámara afectada por un incidente (RF-20).
\end{itemize}

\subsection{Clasificación de riesgo}

Conforme al Reglamento (UE) 2024/1689 (AI Act), ambos componentes se clasifican como riesgo mínimo. Se aplican las obligaciones de transparencia del Artículo 50. Los datos de ventas con nombre/RUC de cliente quedan sujetos a la LOPDP \cite{lopd2021}.

\subsection{Ficha IA-01: Predictor de Demanda por Tipo de Corte}

\begin{longtable}{|p{3cm}|p{11cm}|}
\caption{Ficha del componente IA-01: Predictor de Demanda por Tipo de Corte}\\
\hline
\textbf{Campo} & \textbf{Contenido} \\
\hline
Identificador & IA-01 \\
\hline
Nombre & Predictor de Demanda por Tipo de Corte \\
\hline
Tarea y tipo & Regresión de series temporales: predice la demanda semanal (en libras) por tipo de corte. \\
\hline
Entradas & Historial de ventas últimos 90 días (RF-10, RF-13), día de semana, semana del mes, precio vigente (RF-19). \\
\hline
Salida & Predicción de demanda semanal en libras por tipo de corte, con intervalo de confianza del 80\%. \\
\hline
Consumidor & Propietaria y administrador general (planificación de compras). \\
\hline
Fallback & Si confianza $<$60\% o modelo no disponible: promedio histórico de las últimas 4 semanas con aviso visible al usuario. \\
\hline
Datos de entrenamiento & Origen: RF-10 (mínimo 6 meses, 180+ registros diarios). Sesgo conocido: semanas con cortes de energía (documentado). Base legal: Art. 6 LOPDP -- ejecución contractual. \\
\hline
Clasificación de riesgo & Riesgo mínimo (AI Act Art. 6). Transparencia ante el usuario (Art. 50). \\
\hline
Plan de monitoreo & Indicador: MAE semanal. Umbral de alerta: MAE $>$15 lb/corte. Reentrenamiento: MAE $>$15 lb por 3 semanas consecutivas. \\
\hline
Trazas & RF-10, RF-13, RF-19, RF-20; RNF-11. \\
\hline
\end{longtable}

\subsubsection{Requisitos Funcionales de IA-01}

\begin{description}[nosep]
\item[RF-IA-01.] El sistema debe generar, para la semana inmediata siguiente, una predicción de demanda en libras por tipo de corte, a partir del historial de ventas de los últimos 90 días.
\item[RF-IA-02.] Para cada predicción generada, el sistema debe mostrar el intervalo de confianza del 80\% en libras.
\item[RF-IA-03.] Cuando la confianza interna del modelo sea inferior al 60\% o el modelo no esté disponible, el sistema debe activar el mecanismo de fallback con aviso visible al usuario.
\end{description}

\subsubsection{Requisitos No Funcionales de IA-01}

\begin{description}[nosep]
\item[RNF-IA-01.] \textbf{Exactitud:} MAE $\leq$ 10 libras por tipo de corte sobre el conjunto de prueba, medido antes de cada despliegue.
\item[RNF-IA-02.] \textbf{Tiempo de inferencia:} generación completa en menos de 5 segundos bajo carga normal.
\item[RNF-IA-03.] \textbf{Escalabilidad:} reentrenamiento con hasta 10.000 registros en menos de 30 minutos sin interrumpir el servicio de inferencia.
\item[RNF-IA-04.] \textbf{Equidad -- sesgo por temporada:} diferencia entre MAE en temporada alta y semanas ordinarias $\leq$ 30\%, medida trimestralmente.
\item[RNF-IA-05.] \textbf{Equidad -- sesgo por tipo de corte:} diferencia de MAE relativo entre el corte con mayor y menor precisión $\leq$ 20 puntos porcentuales.
\item[RNF-IA-06.] \textbf{Explicabilidad:} para cada predicción, el sistema debe mostrar los tres factores con mayor peso en la estimación, en lenguaje natural, en $\leq$ 2 segundos.
\end{description}

\subsection{Ficha IA-02: Detector de Anomalías de Cadena de Frío}

\begin{longtable}{|p{3cm}|p{11cm}|}
\caption{Ficha del componente IA-02: Detector de Anomalías de Cadena de Frío}\\
\hline
\textbf{Campo} & \textbf{Contenido} \\
\hline
Identificador & IA-02 \\
\hline
Nombre & Detector de Anomalías de Cadena de Frío \\
\hline
Tarea y tipo & Clasificación binaria: ``riesgo alto'' o ``riesgo bajo'' de pérdida de cadena de frío ante un incidente de corte de energía. \\
\hline
Entradas & Registros de incidentes RF-20 (fecha, hora, duración, cámaras afectadas) y tipo de producto almacenado en cada cámara (RF-07). \\
\hline
Salida & Clasificación ``riesgo alto'' / ``riesgo bajo'', con probabilidad de clase y recomendación de acción en lenguaje natural. \\
\hline
Consumidor & Administrador general (decisión de iniciar baja RF-09). \\
\hline
Fallback & Regla determinística: corte $>$2h para pollo o $>$4h para res/cerdo = ``riesgo alto'' automático, en $\leq$ 500ms. \\
\hline
Datos de entrenamiento & Origen: RF-20 y RF-09 (mínimo 12 meses, 50 incidentes etiquetados). Base legal: no aplica LOPDP (sin datos personales). \\
\hline
Clasificación de riesgo & Riesgo mínimo (AI Act Art. 6). Transparencia ante el usuario (Art. 50). \\
\hline
Plan de monitoreo & Indicador: tasa de falsos negativos mensual. Umbral de alerta: $>$10\% durante el mes. Reentrenamiento: $>$10\% por 2 meses consecutivos o 20 nuevos incidentes etiquetados. \\
\hline
Trazas & RF-09, RF-20, RNF-02 (disponibilidad), RNF-07 (auditoría de bajas). \\
\hline
\end{longtable}

\subsubsection{Requisitos Funcionales de IA-02}

\begin{description}[nosep]
\item[RF-IA-04.] Al registrar un incidente (RF-20), el sistema debe clasificar automáticamente el riesgo como ``alto'' o ``bajo'' y mostrarlo al administrador en la misma pantalla.
\item[RF-IA-05.] Para cada clasificación, el sistema debe mostrar la probabilidad de ``riesgo alto'' y una recomendación de acción en lenguaje natural.
\item[RF-IA-06.] El sistema debe permitir al administrador registrar la acción real tomada y el motivo, para retroalimentar el modelo.
\end{description}

\subsubsection{Requisitos No Funcionales de IA-02}

\begin{description}[nosep]
\item[RNF-IA-07.] \textbf{Sensibilidad:} recall de la clase ``riesgo alto'' $\geq$ 0,90 sobre el conjunto de prueba.
\item[RNF-IA-08.] \textbf{Tiempo de inferencia:} clasificación en menos de 2 segundos desde que el usuario confirma el incidente.
\item[RNF-IA-09.] \textbf{Disponibilidad del fallback:} el mecanismo determinístico debe activarse en $\leq$ 500ms cuando el modelo no esté disponible.
\item[RNF-IA-10.] \textbf{Equidad -- sesgo por tipo de producto:} diferencia de sensibilidad entre la clase ``pollo'' y ``res/cerdo'' $\leq$ 10 puntos porcentuales.
\item[RNF-IA-11.] \textbf{Equidad -- sesgo por horario:} diferencia de tasa de error entre incidentes nocturnos y diurnos $\leq$ 15 puntos porcentuales, medida mensualmente.
\item[RNF-IA-12.] \textbf{Explicabilidad:} para cada clasificación de riesgo alto, el sistema debe mostrar los dos factores más influyentes en lenguaje natural en $\leq$ 2 segundos.
\end{description}


\subsection{Requisitos no funcionales transversales del componente inteligente}

Los requisitos de esta subsección se aplican a los tres componentes
inteligentes del sistema: IA-01 (predictor de demanda), IA-02 (detector de
anomalías de cadena de frío) y DET-01 (detector de patrones de ambigüedad en
requisitos, empleado en el componente empírico). Cada uno declara métrica,
unidad, umbral, método de verificación, responsable y frecuencia de medición,
conforme al §8 de la guía de la asignatura.

% ---------------------------------------------------------------------------
% Macro de ficha. Si el documento ya define una equivalente, reutilicela y
% borre esta definicion para no duplicar.
% ---------------------------------------------------------------------------
\newcommand{\rnfia}[9]{%
\begin{longtable}{|p{3.2cm}|p{10.6cm}|}
\caption[#1: #2]{Requisito no funcional #1 -- #2}\\
\hline
\multicolumn{2}{|l|}{\textbf{#1 -- #2}} \\
\hline
Enunciado verificable & #3 \\
\hline
Métrica y unidad & #4 \\
\hline
Umbral & #5 \\
\hline
Método de verificación & #6 \\
\hline
Responsable & #7 \\
\hline
Frecuencia de medición & #8 \\
\hline
Trazas & #9 \\
\hline
\end{longtable}
\vspace{-0.4cm}
}

\subsubsection{Componente de detección de ambigüedad (DET-01)}

\rnfia{RNF-IA-13}{Exactitud del detector de ambigüedad}
{El detector de patrones de ambigüedad debe alcanzar una exhaustividad no
inferior a 0,70 frente al consenso de un panel de al menos tres personas
expertas que clasifiquen el mismo corpus de forma ciega e independiente.}
{Exhaustividad (\textit{recall}) y medida $F_1$, con intervalo de confianza al
95\,\% por \textit{bootstrap} de percentiles; adimensional.}
{Exhaustividad $\geq$ 0,70 y $F_1 \geq$ 0,60. Si el límite inferior del
intervalo de confianza queda por debajo de 0,50, el resultado se declara no
concluyente y no se emplea como base de decisión.}
{Ejecución de \texttt{07\_Datos/scripts/run\_all.py} sobre el corpus completo
extraído del ERS, contrastada con \texttt{etiquetas\_expertos.csv}. Salida
verificable en \texttt{07\_Datos/resultados/matriz\_confusion\_prf1.json} y
\texttt{bootstrap\_ic95.json}.}
{Verificador de calidad del equipo.}
{En cada versión del ERS y en cada modificación del corpus de requisitos.}
{DET-01; RF-01 a RF-27; caso de prueba CP-IA-13.}

\begin{quote}
\small
\textbf{Estado de cumplimiento en la versión actual.} No se cumple. La
ejecución sobre los 27 RF de la v2.0 arroja exhaustividad $=0{,}0000$ y
$F_1 = 0{,}0000$, con VP $=0$, FP $=0$, FN $=4$ y VN $=23$. El diagnóstico,
verificado patrón por patrón, es que el detector es inerte sobre este corpus:
el umbral de conjunciones múltiples exige más de tres conectores y el máximo
observado es dos; ninguno de los veintiséis patrones de cuantificador vago
aparece; y los veintisiete requisitos emplean la forma activa. Se declara el
incumplimiento en lugar de ajustar el umbral, porque modificar la lógica del
detector después de conocer los resultados invalidaría la comparación
pre-registrada. Con VP $=0$ y FP $=0$ la precisión es una cantidad indefinida,
reportada como cero por convención, y el intervalo de confianza de anchura
nula es un artefacto de remuestrear una constante.
\end{quote}

\subsubsection{Supervisión humana}

\rnfia{RNF-IA-14}{Supervisión humana efectiva de toda marca automática}
{Ninguna salida de un componente inteligente produce efectos sobre el
inventario, la contabilidad o la especificación sin aceptación o descarte
explícito de una persona con rol autorizado. El sistema registra, para cada
salida, la decisión humana adoptada, su autoría, el instante y el motivo
cuando se descarta.}
{Proporción de salidas del componente con decisión humana registrada;
porcentaje. Latencia entre la emisión de la salida y la decisión; horas.}
{100\,\% de las salidas con decisión registrada. Ninguna salida permanece sin
resolver más de 48 horas: transcurrido ese plazo se escala automáticamente al
administrador general.}
{Auditoría del registro de decisiones: consulta que cruce el total de salidas
emitidas con el total de decisiones registradas. La diferencia debe ser cero.
Prueba de aceptación: emitir una salida y comprobar que el efecto sobre el
inventario no se produce hasta que se registra la decisión.}
{Administrador general (aceptación y descarte); verificador de calidad
(auditoría del registro).}
{Auditoría mensual del registro; verificación en cada despliegue.}
{IA-01; IA-02; DET-01; RF-09; RF-IA-06; RNF-07; caso de prueba CP-IA-14.}

\subsubsection{Monitoreo posterior al despliegue}

\rnfia{RNF-IA-15}{Monitoreo continuo en producción}
{Cada componente inteligente en producción se monitoriza de forma continua
mediante indicadores declarados, con periodicidad fija y umbrales de alerta
que disparan revisión o reentrenamiento. El resultado de cada medición queda
registrado y es consultable.}
{IA-01: error absoluto medio semanal, en libras. IA-02: tasa de falsos
negativos mensual, en porcentaje. DET-01: exhaustividad frente a
re-anotación humana periódica, adimensional. Transversal: deriva de la
distribución de entradas respecto de la de entrenamiento, medida con el índice
de estabilidad poblacional (PSI), adimensional.}
{IA-01: alerta si MAE $>$ 15 lb por corte; reentrenamiento si se supera
durante 3 semanas consecutivas. IA-02: alerta si la tasa de falsos negativos
$>$ 10\,\% en el mes; reentrenamiento si se supera 2 meses seguidos o al
acumular 20 incidentes nuevos etiquetados. DET-01: revisión si la
exhaustividad cae más de 0,10 respecto de la medición anterior. Transversal:
alerta si PSI $>$ 0,20.}
{Panel de monitoreo con serie histórica de cada indicador y registro de
alertas disparadas. Revisión documentada de cada alerta, con la acción
adoptada.}
{Administrador general (revisión de alertas); equipo técnico
(instrumentación y panel).}
{IA-01: semanal. IA-02: mensual. DET-01: en cada versión del ERS.
Deriva: mensual. Revisión integral: trimestral.}
{IA-01; IA-02; DET-01; RNF-02; caso de prueba CP-IA-15.}

\subsubsection{Clasificación del nivel de riesgo}

\rnfia{RNF-IA-16}{Clasificación y revisión del nivel de riesgo}
{Todo componente inteligente tiene asignado y documentado un nivel de riesgo,
con la norma de referencia, el criterio aplicado y las obligaciones que se
derivan. La clasificación se revisa cuando cambia el alcance del componente,
su población de usuarios o el uso que se hace de su salida.}
{Nivel de riesgo asignado, en escala ordinal \{mínimo, limitado, alto,
inaceptable\}; antigüedad de la última revisión, en meses.}
{Todo componente en producción con nivel asignado y antigüedad de revisión
no superior a 12 meses. Ningún componente clasificado como inaceptable puede
desplegarse.}
{Revisión documentada frente a los criterios de la norma de referencia, con
firma del responsable y fecha. Registro de la revisión en el repositorio.}
{Líder del equipo, con visto bueno del docente supervisor.}
{Anual, y de forma extraordinaria ante cualquier cambio de alcance o de uso.}
{IA-01; IA-02; DET-01; caso de prueba CP-IA-16.}

\begin{quote}
\small
\textbf{Clasificación vigente.} Los tres componentes se clasifican como de
\textbf{riesgo mínimo} conforme al Reglamento (UE) 2024/1689. Ninguno figura
entre los sistemas prohibidos del Artículo 5 ni entre los de alto riesgo del
Artículo 6 y su Anexo III: no realizan identificación biométrica, no evalúan a
personas ni condicionan su acceso a servicios, y su salida está sujeta en todo
caso a decisión humana (RNF-IA-14). Se aplican, por tanto, las obligaciones de
transparencia del Artículo 50: la persona usuaria es informada de que
interactúa con un sistema de inteligencia artificial y de cuál es el alcance de
su salida.

Se corrige aquí una imprecisión de la versión anterior del documento, que
atribuía la categoría de riesgo mínimo al Artículo 6. El Artículo 6 regula la
clasificación de los sistemas de \emph{alto} riesgo; el riesgo mínimo es la
categoría residual que resulta de no encajar en los Artículos 5 ni 6. El
tratamiento de datos personales asociado se rige por la Ley Orgánica de
Protección de Datos Personales del Ecuador.
\end{quote}

\subsubsection{Explicabilidad y equidad: consolidación}

Los requisitos de explicabilidad (RNF-IA-06, RNF-IA-12) y de equidad
(RNF-IA-04, RNF-IA-05, RNF-IA-10, RNF-IA-11), ya especificados en las fichas de
IA-01 e IA-02, se extienden al detector DET-01 mediante los dos requisitos
siguientes, y se completan con el responsable y la frecuencia de medición que
la guía exige y que aquellos no declaraban.

\rnfia{RNF-IA-17}{Explicabilidad del detector de ambigüedad}
{Para cada marca de ambigüedad, el sistema muestra al analista la categoría de
patrón activada, el fragmento literal del requisito que la activó y la regla
concreta que se aplicó, sin que el analista deba consultar el código fuente.}
{Proporción de marcas acompañadas de evidencia literal, en porcentaje; tiempo
de presentación de la explicación, en segundos.}
{100\,\% de las marcas con evidencia literal; presentación en $\leq$ 2
segundos.}
{Inspección de \texttt{clasificaciones\_detector.csv}: las columnas
\texttt{evidencia\_*} deben estar cumplimentadas en toda fila con
\texttt{ambiguo\_detector} $=1$. Prueba de aceptación sobre un requisito
sintético que active cada categoría.}
{Documentador del equipo.}
{En cada ejecución del pipeline de análisis.}
{DET-01; RNF-IA-13; caso de prueba CP-IA-17.}

\rnfia{RNF-IA-18}{Equidad del detector entre categorías de requisito}
{El detector no debe presentar diferencias sistemáticas de exhaustividad entre
los distintos tipos de requisito funcional del corpus, de modo que ninguna
familia de requisitos quede sistemáticamente sin revisar.}
{Diferencia de exhaustividad entre el grupo de requisitos con mayor y con menor
detección, agrupados por módulo funcional; puntos porcentuales.}
{Diferencia $\leq$ 20 puntos porcentuales. Cuando el detector no marque ningún
requisito, la métrica se declara no evaluable y así se reporta, en lugar de
computarse como diferencia nula.}
{Cálculo por script sobre \texttt{dataset\_consolidado.csv}, desagregando por
módulo funcional según la matriz de trazabilidad.}
{Verificador de calidad del equipo.}
{En cada versión del ERS.}
{DET-01; RNF-IA-13; caso de prueba CP-IA-18.}


%==========================================================
\section{Auditoría de Calidad del ERS con Seis Métricas}
%==========================================================

La auditoría se aplica sobre el universo de 42 requisitos base (27 RF + 15 RNF) del ERS/SRS v2.0, siguiendo las métricas derivadas de ISO/IEC 25010:2023 \cite{iso25010}.

\subsection{Instrumento de auditoría -- Conteos base}

\begin{itemize}[nosep]
\item Total de requisitos base: 27 RF + 15 RNF = \textbf{42 requisitos}.
\item Requisitos con criterio BDD comprobable (v1.0/Entrega 2A): 37/40.
\item Requisitos con criterio BDD comprobable (v2.0/Entrega 4 2B): 41/42.
\item Casos de uso identificados: 12; CU especificados: 12.
\item Actores: 5; actores con $\geq$1 RF: 5.
\item Pares de requisitos analizados: $\binom{42}{2} = 861$ (análisis exhaustivo).
\item Conflictos detectados (v1.0): 1 (D-03).
\item Defectos residuales en re-inspección PE5 (v2.0): 0.
\end{itemize}

\subsection{Cálculo de las seis métricas}

\subsubsection{M1 -- Completitud}

\begin{align*}
M1a_{\text{antes}} &= \frac{38}{40} \times 100 = 95{,}0\% \\
M1a_{\text{después}} &= \frac{41}{42} \times 100 = 97{,}6\% \\
M1b &= \frac{12}{12} \times 100 = 100{,}0\% \\
M1c &= \frac{5}{5} \times 100 = 100{,}0\%
\end{align*}

\textbf{Veredicto M1:} $\checkmark$ CUMPLE (M1a $\geq$ 95\%, M1b = 100\%, M1c = 100\%). RF-18 (Could) se mantiene sin criterio BDD formal, documentado como excepción justificada por prioridad baja.

\subsubsection{M2 -- Consistencia}

\begin{align*}
M2_{\text{antes}} &= 1 - \frac{1}{780} = 0{,}9987 \\
M2_{\text{después}} &= 1 - \frac{0}{861} = 1{,}0000
\end{align*}

\textbf{Veredicto M2:} $\checkmark$ CUMPLE (M2 $\geq$ 0,98; cero conflictos abiertos en v2.0).

\subsubsection{M3 -- Verificabilidad}

\begin{align*}
M3_{\text{antes}} &= \frac{37}{40} \times 100 = 92{,}5\% \\
M3_{\text{después}} &= \frac{41}{42} \times 100 = 97{,}6\%
\end{align*}

\textbf{Veredicto M3:} $\checkmark$ CUMPLE (M3 $\geq$ 90\%; RF críticos al 100\%).

\subsubsection{M4 -- Trazabilidad}

\begin{align*}
M4_{ade\text{ antes}} &= \frac{23}{25} \times 100 = 92{,}0\% \\
M4_{ade\text{ después}} &= \frac{26}{27} \times 100 = 96{,}3\% \\
M4_{atr} &= \frac{27}{27} \times 100 = 100{,}0\%
\end{align*}

\textbf{Veredicto M4:} $\checkmark$ CUMPLE (M4ade $\geq$ 90\%; M4atr = 100\%). RF-18 (Could) sin HU, justificado por prioridad baja.

\subsubsection{M5 -- Modificabilidad}

\begin{longtable}{|p{1.5cm}|p{4cm}|p{7cm}|p{1.5cm}|}
\caption{Cálculo de acoplamiento para M5 (Modificabilidad)}\\
\hline
\textbf{RF} & \textbf{Razón de selección} & \textbf{RF afectados al modificar} & \textbf{N.° afect.} \\
\hline
RF-03 & Nodo de mayor dependencia & RF-04, RF-06, RF-07, RF-11 & 4 \\
RF-07 & Hub de inventario & RF-08, RF-09, RF-10, RF-12, RF-16, RF-24 & 6 \\
RF-10 & Ventas (salida principal) & RF-13, RF-14, RF-22, RF-27 & 4 \\
RF-04 & Pesaje (cálculo central) & RF-05, RF-19 & 2 \\
RF-11 & Trazabilidad (consulta) & RF-23, RF-25, RF-26 & 3 \\
RF-15 & Gestión de usuarios & RF-01 & 1 \\
RF-20 & Incidentes de energía & RF-09, RF-08 & 2 \\
\hline
\textbf{Total} & & & \textbf{22} \\
\hline
\end{longtable}

$$M5 = \frac{4 + 6 + 4 + 2 + 3 + 1 + 2}{7} = \frac{22}{7} = 3{,}14 \text{ requisitos afectados en promedio}$$

\textbf{Nota:} El valor M5 = 3,14 supera levemente el umbral de referencia de 3,0 por la incorporación de RF-26 y RF-27 (que amplían las dependencias de RF-11 y RF-10 respectivamente). El equipo documentó esta situación: la arquitectura en capas (RNF-05) y el patrón de repositorio del módulo de trazabilidad mitigan el impacto real de modificación, dado que RF-26 comparte el componente \texttt{Backend-Trazabilidad} con RF-11 y RF-25 sin requerir cambios en la capa de datos.

\textbf{Veredicto M5:} $\approx$ ACEPTABLE (3,14; desviación mínima del umbral documentada y mitigada por RNF-05).

\subsubsection{M6 -- Corrección}

\begin{align*}
M6_{\text{antes}} &= \frac{4}{40} = 0{,}10 \text{ defectos por requisito} \\
M6_{\text{después}} &= \frac{0}{42} = 0{,}00 \text{ defectos por requisito}
\end{align*}

\textbf{Veredicto M6:} $\checkmark$ CUMPLE (M6 $\leq$ 0,05; todos los defectos de re-inspección corregidos).

\subsection{Tabla consolidada de auditoría (antes/después)}

\begin{longtable}{|p{1cm}|p{3cm}|p{1.8cm}|p{2cm}|p{2cm}|p{2.5cm}|p{1cm}|}
\caption{Tabla consolidada de auditoría de calidad ERS/SRS: antes (v1.0) y después (v2.0)}\\
\hline
\textbf{M.} & \textbf{Métrica} & \textbf{Antes v1.0} & \textbf{Después v2.0} & \textbf{Valor ref.} & \textbf{Ref. normativa} & \textbf{Cumple} \\
\hline
M1a & Completitud atributos & 95,0\% & 97,6\% & $\geq$ 95\% & ISO/IEC/IEEE 29148:2018 & Sí \\
\hline
M1b & CU especificados & 100,0\% & 100,0\% & 100\% & ISO/IEC/IEEE 29148:2018 & Sí \\
\hline
M1c & Actores con RF & 100,0\% & 100,0\% & 100\% & ISO/IEC/IEEE 29148:2018 & Sí \\
\hline
M2 & Consistencia & 0,9987 & 1,0000 & $\geq$ 0,98 & ISO/IEC 25010:2023 & Sí \\
\hline
M3 & Verificabilidad & 92,5\% & 97,6\% & $\geq$ 90\% & ISO/IEC/IEEE 29148:2018 & Sí \\
\hline
M4ade & Traza adelante & 92,0\% & 96,3\% & $\geq$ 90\% & ISO/IEC/IEEE 29148:2018 & Sí \\
\hline
M4atr & Traza hacia atrás & 100,0\% & 100,0\% & 100\% & ISO/IEC/IEEE 29148:2018 & Sí \\
\hline
M5 & Modificabilidad & 2,71 req & 3,14 req & $\leq$ 3,0 & ISO/IEC 25010:2023 & Aprox. \\
\hline
M6 & Corrección & 0,10 def/req & 0,00 def/req & $\leq$ 0,05 & ISO/IEC 25010:2023 & Sí \\
\hline
\end{longtable}

\textbf{Resultado global: 5/6 métricas cumplen el valor de referencia; M5 presenta desviación mínima documentada y mitigada.}

%==========================================================
% APÉNDICES
%==========================================================
\appendix

\section{Plan del proyecto de Ingeniería de Requisitos}

\subsection{Cronograma de actividades de IR}

\begin{longtable}{|p{1.6cm}|p{8.5cm}|p{3.5cm}|}
\caption{Cronograma de actividades de ingeniería de requisitos}\\
\hline
\textbf{Semana} & \textbf{Actividad} & \textbf{Responsable} \\
\hline
1--10 & Actividades de la Entrega 1A y 1B (elicitación inicial, mockups, RF/RNF parciales, modelado UML parcial, MoSCoW). & Todo el equipo \\
\hline
11 & Adenda ética de la segunda ronda; diseño experimental; registro del protocolo en OSF. & Verificador / Analista líder \\
\hline
12 & Ejecución de la segunda ronda de campo; cuestionario; validación piloto. & Todo el equipo \\
\hline
13 & Consolidación del ERS/SRS completo (Entrega 2A / 3) y del MVP. & Todo el equipo \\
\hline
14 & Selección de revista objetivo; cierre del 80\% de datos primarios. & Todo el equipo \\
\hline
15 & Análisis empírico completo; borrador de resultados; sesión member-checking. & Todo el equipo \\
\hline
16 & Cierre del manuscrito; depósito Zenodo; autoevaluación FAIR; simulacro de defensa. & Todo el equipo \\
\hline
\textbf{17} & \textbf{Entrega 4 (2B) y defensa final.} Corte: domingo 23/08/2026, 23:59. & \textbf{Todo el equipo} \\
\hline
\end{longtable}

%==========================================================
\section{Evidencias de elicitación}
%==========================================================

\subsection{Instrumento de recolección de datos (guía de entrevista)}

\begin{figure}[H]
    \centering
    \includegraphics[width=0.9\textwidth,keepaspectratio]{figures/A02_Instrumentos_Recoleccion_page-0001.jpg}
    \caption{Guía de entrevista -- Página 1}
\end{figure}
\begin{figure}[H]
    \centering
    \includegraphics[width=0.9\textwidth,keepaspectratio]{figures/A02_Instrumentos_Recoleccion_page-0002.jpg}
    \caption{Guía de entrevista -- Página 2}
\end{figure}
\begin{figure}[H]
    \centering
    \includegraphics[width=0.9\textwidth,keepaspectratio]{figures/A02_Instrumentos_Recoleccion_page-0003.jpg}
    \caption{Guía de entrevista -- Página 3}
\end{figure}
\begin{figure}[H]
    \centering
    \includegraphics[width=0.9\textwidth,keepaspectratio]{figures/A02_Instrumentos_Recoleccion_page-0004.jpg}
    \caption{Guía de entrevista -- Página 4}
\end{figure}
\begin{figure}[H]
    \centering
    \includegraphics[width=0.9\textwidth,keepaspectratio]{figures/A02_Instrumentos_Recoleccion_page-0005.jpg}
    \caption{Guía de entrevista -- Página 5}
\end{figure}

\subsection{Consentimientos informados}

\begin{figure}[H]
    \centering
    \includegraphics[width=0.9\textwidth,keepaspectratio]{figures/2026-07-28_Propietaria_ENTR-05_Consentimiento_page-0001.jpg}
    \caption{Consentimiento informado -- Propietaria (ENTR-05), página 1}
\end{figure}
\begin{figure}[H]
    \centering
    \includegraphics[width=0.9\textwidth,keepaspectratio]{figures/2026-07-28_Propietaria_ENTR-05_Consentimiento_page-0002.jpg}
    \caption{Consentimiento informado -- Propietaria (ENTR-05), página 2}
\end{figure}
\begin{figure}[H]
    \centering
    \includegraphics[width=0.9\textwidth,keepaspectratio]{figures/2026-08-01_AdministradorDeBodega_ENTR-01_Consentimiento_page-0001.jpg}
    \caption{Consentimiento informado -- Administrador de Bodega (ENTR-01), página 1}
\end{figure}
\begin{figure}[H]
    \centering
    \includegraphics[width=0.9\textwidth,keepaspectratio]{figures/2026-08-01_AdministradorDeBodega_ENTR-01_Consentimiento_page-0002.jpg}
    \caption{Consentimiento informado -- Administrador de Bodega (ENTR-01), página 2}
\end{figure}
\begin{figure}[H]
    \centering
    \includegraphics[width=0.9\textwidth,keepaspectratio]{figures/2026-08-01_AdministradorDeBodega_ENTR-02_Consentimiento_page-0001.jpg}
    \caption{Consentimiento informado -- Administrador de Bodega (ENTR-02), página 1}
\end{figure}
\begin{figure}[H]
    \centering
    \includegraphics[width=0.9\textwidth,keepaspectratio]{figures/2026-08-01_AdministradorDeBodega_ENTR-02_Consentimiento_page-0002.jpg}
    \caption{Consentimiento informado -- Administrador de Bodega (ENTR-02), página 2}
\end{figure}
\begin{figure}[H]
    \centering
    \includegraphics[width=0.9\textwidth,keepaspectratio]{figures/AdminBodega1.jpg}
    \caption{Evidencia fotográfica -- Administrador de Bodega, sesión 1}
\end{figure}
\begin{figure}[H]
    \centering
    \includegraphics[width=0.9\textwidth,keepaspectratio]{figures/AdminBodega2.jpg}
    \caption{Evidencia fotográfica -- Administrador de Bodega, sesión 2}
\end{figure}
\begin{figure}[H]
    \centering
    \includegraphics[width=0.9\textwidth,keepaspectratio]{figures/2026-08-01_CarniceroPrincipal_ENTR-03_Consentimiento_page-0001.jpg}
    \caption{Consentimiento informado -- Carnicero Principal (ENTR-03), página 1}
\end{figure}
\begin{figure}[H]
    \centering
    \includegraphics[width=0.9\textwidth,keepaspectratio]{figures/2026-08-01_CarniceroPrincipal_ENTR-03_Consentimiento_page-0002.jpg}
    \caption{Consentimiento informado -- Carnicero Principal (ENTR-03), página 2}
\end{figure}
\begin{figure}[H]
    \centering
    \includegraphics[width=0.9\textwidth,keepaspectratio]{figures/2026-08-01_OperarioRecepcionPesaje_ENTR-04_Consentimiento_page-0001.jpg}
    \caption{Consentimiento informado -- Operario Recepción/Pesaje (ENTR-04), página 1}
\end{figure}
\begin{figure}[H]
    \centering
    \includegraphics[width=0.9\textwidth,keepaspectratio]{figures/2026-08-01_OperarioRecepcionPesaje_ENTR-04_Consentimiento_page-0002.jpg}
    \caption{Consentimiento informado -- Operario Recepción/Pesaje (ENTR-04), página 2}
\end{figure}
\begin{figure}[H]
    \centering
    \includegraphics[width=0.9\textwidth,keepaspectratio]{figures/2026-08-29_Chofer_ENTR-06_Consentimiento_page-0001.jpg}
    \caption{Consentimiento informado -- Chofer (ENTR-06), página 1}
\end{figure}
\begin{figure}[H]
    \centering
    \includegraphics[width=0.9\textwidth,keepaspectratio]{figures/2026-08-29_Chofer_ENTR-06_Consentimiento_page-0002.jpg}
    \caption{Consentimiento informado -- Chofer (ENTR-06), página 2}
\end{figure}
\begin{figure}[H]
    \centering
    \includegraphics[width=0.9\textwidth,keepaspectratio]{figures/2026-08-30_Recepcion_Bodega_ENTR-07_Consentimiento_page-0001.jpg}
    \caption{Consentimiento informado -- Recepción de Bodega (ENTR-07), página 1}
\end{figure}
\begin{figure}[H]
    \centering
    \includegraphics[width=0.9\textwidth,keepaspectratio]{figures/2026-08-30_Recepcion_Bodega_ENTR-07_Consentimiento_page-0002.jpg}
    \caption{Consentimiento informado -- Recepción de Bodega (ENTR-07), página 2}
\end{figure}
\begin{figure}[H]
    \centering
    \includegraphics[width=0.9\textwidth,keepaspectratio]{figures/2026-08-31_Operario_RecepcionPesaje_ENTR-08_Consentimiento_page-0001.jpg}
    \caption{Consentimiento informado -- Operario Recepción/Pesaje (ENTR-08), página 1}
\end{figure}
\begin{figure}[H]
    \centering
    \includegraphics[width=0.9\textwidth,keepaspectratio]{figures/2026-08-31_Operario_RecepcionPesaje_ENTR-08_Consentimiento_page-0002.jpg}
    \caption{Consentimiento informado -- Operario Recepción/Pesaje (ENTR-08), página 2}
\end{figure}
\begin{figure}[H]
    \centering
    \includegraphics[width=0.9\textwidth,keepaspectratio]{figures/2026-08-31_Vendedor_Mostrador_ENTR-09_Consentimiento_page-0001.jpg}
    \caption{Consentimiento informado -- Vendedor de Mostrador (ENTR-09), página 1}
\end{figure}
\begin{figure}[H]
    \centering
    \includegraphics[width=0.9\textwidth,keepaspectratio]{figures/2026-08-31_Vendedor_Mostrador_ENTR-09_Consentimiento_page-0002.jpg}
    \caption{Consentimiento informado -- Vendedor de Mostrador (ENTR-09), página 2}
\end{figure}
\begin{figure}[H]
    \centering
    \includegraphics[width=0.9\textwidth,keepaspectratio]{figures/2026-08-31_Vendedor_Mostrador_ENTR-10_Consentimiento_page-0001.jpg}
    \caption{Consentimiento informado -- Vendedor de Mostrador (ENTR-10), página 1}
\end{figure}
\begin{figure}[H]
    \centering
    \includegraphics[width=0.9\textwidth,keepaspectratio]{figures/2026-08-31_Vendedor_Mostrador_ENTR-10_Consentimiento_page-0002.jpg}
    \caption{Consentimiento informado -- Vendedor de Mostrador (ENTR-10), página 2}
\end{figure}
\begin{figure}[H]
    \centering
    \includegraphics[width=0.9\textwidth,keepaspectratio]{figures/2026-09-01_Encargada_Sucursal_ENTR-11_Consentimiento_page-0001.jpg}
    \caption{Consentimiento informado -- Encargada de Sucursal (ENTR-11), página 1}
\end{figure}
\begin{figure}[H]
    \centering
    \includegraphics[width=0.9\textwidth,keepaspectratio]{figures/2026-09-01_Encargada_Sucursal_ENTR-11_Consentimiento_page-0002.jpg}
    \caption{Consentimiento informado -- Encargada de Sucursal (ENTR-11), página 2}
\end{figure}
\begin{figure}[H]
    \centering
    \includegraphics[width=0.9\textwidth,keepaspectratio]{figures/2026-09-01_Administrador_Contador_ENTR-12_Consentimiento_page-0001.jpg}
    \caption{Consentimiento informado -- Administrador/Contador (ENTR-12), página 1}
\end{figure}
\begin{figure}[H]
    \centering
    \includegraphics[width=0.9\textwidth,keepaspectratio]{figures/2026-09-01_Administrador_Contador_ENTR-12_Consentimiento_page-0002.jpg}
    \caption{Consentimiento informado -- Administrador/Contador (ENTR-12), página 2}
\end{figure}
\begin{figure}[H]
    \centering
    \includegraphics[width=0.9\textwidth,keepaspectratio]{figures/2026-09-01_Chofer_Despacho_ENTR-13_Consentimiento_page-0001.jpg}
    \caption{Consentimiento informado -- Chofer de Despacho (ENTR-13), página 1}
\end{figure}
\begin{figure}[H]
    \centering
    \includegraphics[width=0.9\textwidth,keepaspectratio]{figures/2026-09-01_Chofer_Despacho_ENTR-13_Consentimiento_page-0002.jpg}
    \caption{Consentimiento informado -- Chofer de Despacho (ENTR-13), página 2}
\end{figure}
\begin{figure}[H]
    \centering
    \includegraphics[width=0.9\textwidth,keepaspectratio]{figures/2026-08-31_Carnicero_Despachador_ENTR-14_Consentimiento_page-0001.jpg}
    \caption{Consentimiento informado -- Carnicero Despachador (ENTR-14), página 1}
\end{figure}
\begin{figure}[H]
    \centering
    \includegraphics[width=0.9\textwidth,keepaspectratio]{figures/2026-08-31_Carnicero_Despachador_ENTR-14_Consentimiento_page-0002.jpg}
    \caption{Consentimiento informado -- Carnicero Despachador (ENTR-14), página 2}
\end{figure}
\begin{figure}[H]
    \centering
    \includegraphics[width=0.9\textwidth,keepaspectratio]{figures/2026-08-31_Chofer_Despachador_ENTR-15_Consentimiento_page-0001.jpg}
    \caption{Consentimiento informado -- Chofer Despachador (ENTR-15), página 1}
\end{figure}
\begin{figure}[H]
    \centering
    \includegraphics[width=0.9\textwidth,keepaspectratio]{figures/2026-08-31_Chofer_Despachador_ENTR-15_Consentimiento_page-0002.jpg}
    \caption{Consentimiento informado -- Chofer Despachador (ENTR-15), página 2}
\end{figure}
\begin{figure}[H]
    \centering
    \includegraphics[width=0.9\textwidth,keepaspectratio]{figures/2026-08-31_Vendedor_Demostrador_ENTR-16_Consentimiento_page-0001.jpg}
    \caption{Consentimiento informado -- Vendedor Demostrador (ENTR-16), página 1}
\end{figure}
\begin{figure}[H]
    \centering
    \includegraphics[width=0.9\textwidth,keepaspectratio]{figures/2026-08-31_Vendedor_Demostrador_ENTR-16_Consentimiento_page-0002.jpg}
    \caption{Consentimiento informado -- Vendedor Demostrador (ENTR-16), página 2}
\end{figure}

\subsection{Actas de entrevista}

\begin{figure}[H]
    \centering
    \includegraphics[width=0.9\textwidth,keepaspectratio]{figures/ACTA-01.jpg.jpg}
    \caption{Acta de entrevista 01}
\end{figure}
\begin{figure}[H]
    \centering
    \includegraphics[width=0.9\textwidth,keepaspectratio]{figures/ACTA-02.jpg}
    \caption{Acta de entrevista 02}
\end{figure}
\begin{figure}[H]
    \centering
    \includegraphics[width=0.9\textwidth,keepaspectratio]{figures/ACTA-03.jpg}
    \caption{Acta de entrevista 03}
\end{figure}
\begin{figure}[H]
    \centering
    \includegraphics[width=0.9\textwidth,keepaspectratio]{figures/ACTA-04.jpg}
    \caption{Acta de entrevista 04}
\end{figure}
\begin{figure}[H]
    \centering
    \includegraphics[width=0.9\textwidth,keepaspectratio]{figures/ACTA-05.jpg}
    \caption{Acta de entrevista 05}
\end{figure}
\begin{figure}[H]
    \centering
    \includegraphics[width=0.9\textwidth,keepaspectratio]{figures/ACTA-06.jpg}
    \caption{Acta de entrevista 06}
\end{figure}
\begin{figure}[H]
    \centering
    \includegraphics[width=0.9\textwidth,keepaspectratio]{figures/ACTA-07.jpg}
    \caption{Acta de entrevista 07}
\end{figure}
\begin{figure}[H]
    \centering
    \includegraphics[width=0.9\textwidth,keepaspectratio]{figures/ACTA-08.jpg}
    \caption{Acta de entrevista 08}
\end{figure}
\begin{figure}[H]
    \centering
    \includegraphics[width=0.9\textwidth,keepaspectratio]{figures/ACTA-09.jpg}
    \caption{Acta de entrevista 09}
\end{figure}
\begin{figure}[H]
    \centering
    \includegraphics[width=0.9\textwidth,keepaspectratio]{figures/ACTA-10.jpg}
    \caption{Acta de entrevista 10}
\end{figure}
\begin{figure}[H]
    \centering
    \includegraphics[width=0.9\textwidth,keepaspectratio]{figures/ACTA-11.jpg}
    \caption{Acta de entrevista 11}
\end{figure}
\begin{figure}[H]
    \centering
    \includegraphics[width=0.9\textwidth,keepaspectratio]{figures/ACTA-12.jpg}
    \caption{Acta de entrevista 12}
\end{figure}
\begin{figure}[H]
    \centering
    \includegraphics[width=0.9\textwidth,keepaspectratio]{figures/ACTA-13.jpg}
    \caption{Acta de entrevista 13}
\end{figure}
\begin{figure}[H]
    \centering
    \includegraphics[width=0.9\textwidth,keepaspectratio]{figures/ACTA-14.jpg}
    \caption{Acta de entrevista 14}
\end{figure}
\begin{figure}[H]
    \centering
    \includegraphics[width=0.9\textwidth,keepaspectratio]{figures/ACTA-15.jpg}
    \caption{Acta de entrevista 15}
\end{figure}
\begin{figure}[H]
    \centering
    \includegraphics[width=0.9\textwidth,keepaspectratio]{figures/ACTA-16.jpg}
    \caption{Acta de entrevista 16}
\end{figure}

%==========================================================
\section{Gestión de cambios: CCB y RFC}
%==========================================================

\subsection{RFC-01 -- Filtro de reportes por proveedor: APROBADO}

El administrador general solicitó que los reportes pudieran filtrarse por proveedor. El CCB aprobó la incorporación de RF-27. Impacto: añade una condición de filtro al componente Backend-Reportes; no modifica el modelo de datos.

\subsection{RFC-02 -- Autenticación multifactor (MFA): APROBADO CON CONDICIÓN}

Se aprobó la implementación de MFA mediante TOTP para roles con capacidad de aprobar bajas y configurar precios (RNF-03). Condición: implementar pregunta de seguridad como fallback cuando no haya conectividad (RF-21 offline).

\subsection{RFC-03 -- Firma veterinaria en exportación de evidencia: APROBADO}

Se aprobó por riesgo regulatorio. Incorporado en RF-02 (registro de guía de origen con firma) y RF-23 (exportación de trazabilidad para Agrocalidad).

%==========================================================
\section{Inspección Fagan PE5 -- Re-inspección y defectos adicionales}
%==========================================================

En la re-inspección PE5 del ERS v1.0 (Entrega 3, 2A) se detectaron 4 defectos adicionales, todos corregidos en v2.0:

\begin{longtable}{|p{1cm}|p{2.5cm}|p{4cm}|p{1.5cm}|p{5cm}|}
\caption{Defectos detectados en la re-inspección Fagan PE5 del ERS v1.0}\\
\hline
\textbf{DEF} & \textbf{Tipo} & \textbf{Descripción} & \textbf{RF/RNF} & \textbf{Acción tomada en v2.0} \\
\hline
D-PE5-01 & Verificabilidad & RNF-05 redactado sin criterio objetivo cuantificable & RNF-05 & Reescrito con umbral cuantificable: ``sin modificar más del 10\% de los módulos adyacentes'' \\
\hline
D-PE5-02 & Completitud & RF-17 sin caso de uso asociado (requisito huérfano) & RF-17 & Enlazado a CU-01 como flujo alternativo \\
\hline
D-PE5-03 & Consistencia & ``Canal'' usada con dos significados distintos en el glosario & RF-02, RF-06 & Definición única añadida al glosario (Sección 1.3): canal = cuerpo del animal faenado sin cabeza ni vísceras \\
\hline
D-PE5-04 & Verificabilidad & RF-08 usaba ``próximo a vencer'' sin umbral temporal definido & RF-08 & Umbral fijado en $\leq$ 1 día; umbrales de vida útil por especie añadidos (pollo: 3d, res: 7d, cerdo: 5d, embutidos: 15d) \\
\hline
\end{longtable}

%==========================================================
\section{Retrospectiva del equipo (Start-Stop-Continue)}
%==========================================================

\begin{longtable}{|p{2cm}|p{6cm}|p{5cm}|}
\caption{Retrospectiva del equipo CarniCore -- Start-Stop-Continue}\\
\hline
\textbf{Categoría} & \textbf{Elemento} & \textbf{Responsable / Horizonte} \\
\hline
\textbf{Start} & Iniciar el registro formal de evidencias (EV-XX) desde la primera entrevista, no retroactivamente. & Castro Bajaña -- Próximo proyecto \\
\hline
\textbf{Start} & Iniciar la revisión de trazabilidad en paralelo con la especificación de requisitos (PE2), no solo al cierre. & Crespo Espinoza -- Próximo proyecto \\
\hline
\textbf{Start} & Iniciar el uso de etiquetas de versión en Git desde la PE1. & Gamarra Araujo -- Próximo proyecto \\
\hline
\textbf{Start} & Iniciar reuniones de sincronización semanales breves (15 min) para detectar dependencias entre artefactos. & Pérez Ruiz -- Próximo proyecto \\
\hline
\textbf{Stop} & Dejar de agregar requisitos sin verificar su impacto en la matriz de trazabilidad. & Pérez Ruiz -- Inmediatamente \\
\hline
\textbf{Stop} & Dejar de redactar RNF en términos cualitativos sin umbral cuantificable. & Todo el equipo -- Inmediatamente \\
\hline
\textbf{Stop} & Dejar de concentrar los commits en un solo integrante. & Castro / Gamarra -- Inmediatamente \\
\hline
\textbf{Continue} & Continuar con el formato Connextra + Gherkin para historias de usuario. & Todo el equipo \\
\hline
\textbf{Continue} & Continuar con la entrevista semiestructurada como técnica principal de elicitación. & Todo el equipo \\
\hline
\textbf{Continue} & Continuar con la priorización combinada MoSCoW + Kano + WSJF. & Todo el equipo \\
\hline
\end{longtable}

%==========================================================
\section{Declaración individual de aporte}
%==========================================================

\begin{longtable}{|p{3cm}|p{7cm}|p{4cm}|}
\caption{Declaración individual de aporte -- Entrega 4 (2B)}\\
\hline
\textbf{Integrante} & \textbf{Aportes principales en Entrega 4 (2B)} & \textbf{Evidencia Git} \\
\hline
Castro Bajaña Ariel Omar & Integración de Introducción y Metodología; actualización i* SR; consistencia editorial del ERS v2.0. & Commits Sec. 1, 2, istar\_SR \\
\hline
Crespo Espinoza Kleber Obed & Revisión UML v2.0; verificación ERS final; control de figuras/tablas; RNF actualizados. & Commits UML, RNF, figuras \\
\hline
Gamarra Araujo Edhu Xavier & Trazabilidad end-to-end 84 filas (60 base/IA + 24 de flujos); estructura del repositorio; control de línea base; CHANGELOG. & Commits trazabilidad, repo \\
\hline
Pérez Ruiz Carlos Andrés & Requisitos de IA (Sec.~8); auditoría de calidad M1--M6 (Sec.~9); conclusiones integradoras. & Commits IA, métricas \\
\hline
Quintero Gende Erick Jahir & Re-inspección Fagan PE5; gestión de RFC/CCB; RNF verificabilidad; retrospectiva; banco de preguntas. & Commits Fagan, RFC, retro \\
\hline
\end{longtable}

%==========================================================
\section{Declaración de uso de Inteligencia Artificial}
%==========================================================

\begin{longtable}{|p{3cm}|p{3cm}|p{4cm}|p{4cm}|}
\caption{Declaración de uso de inteligencia artificial -- CarniCore PE5}\\
\hline
\textbf{Sección del ERS} & \textbf{Herramienta} & \textbf{Tipo de asistencia} & \textbf{Método de validación} \\
\hline
Secciones 1--5 (Intro, Descripción, RF, UML, Trazabilidad) & Claude (Anthropic) y ChatGPT (OpenAI) & Corrección de estilo y ortografía; mejora de coherencia entre párrafos & Revisión íntegra por todos los integrantes contra los artefactos del ERS v2.0; ninguna afirmación técnica fue aceptada sin contraste con el documento fuente \\
\hline
Secciones 6--7 (MVP, Empírico) & Claude (Anthropic) & Revisión de redacción técnica; formateo de tablas LaTeX & Verificación contra el repositorio GitHub y el registro OSF; todos los identificadores de evidencia cotejados uno a uno \\
\hline
Sección 8 (Requisitos de IA) & Claude (Anthropic); Google Scholar para referencias & Sugerencia de estructura de fichas IA-01/IA-02; corrección de RNF de equidad y explicabilidad & DOI de cada referencia verificado en \url{https://doi.org}; umbrales validados por el equipo contra los datos del dominio \\
\hline
Sección 9 (Métricas) & Calculadora científica y Google Sheets & Verificación aritmética independiente de M1--M6 & Conteos base realizados manualmente sobre el ERS v2.0 por dos integrantes de forma independiente \\
\hline
Apéndices (Retro, Conclusiones) & Claude (Anthropic) & Revisión de ortografía y cohesión & Contenido redactado y validado íntegramente por el equipo; la IA no generó juicios ni análisis \\
\hline
\end{longtable}

Las secciones evaluativas (análisis, justificaciones de decisiones de IR, conclusiones integradoras) son producción propia del equipo, verificada por todos los integrantes. El uso de asistentes de IA se limitó a corrección de estilo, ortografía y formato.

\newpage
\bibliographystyle{plain}
\bibliography{referencias}

\end{document}
